\chapter{Absolute Hodge Classes, the Tate and Standard Conjectures}\label{absolute-hodge}
\chapterstatus{complete}

\section{Introduction}\label{absolute-hodge:section-introduction}

The cohomology of a nonsingular projective variety over $\CC$ comes in several flavours: Betti cohomology with its rational structure, algebraic de Rham cohomology with its Hodge filtration, and
$\ell$-adic cohomology with, for varieties defined over smaller fields, a Galois action. The Hodge conjecture concerns classes that are special in the first two at once. An automorphism
$\sigma$ of the field $\CC$, in general wildly discontinuous, transforms a variety $X$ into a \emph{conjugate variety} $X^\sigma$ and acts on de Rham and $\ell$-adic cohomology, but not on Betti
cohomology. Classes of algebraic cycles are carried to classes of algebraic cycles, so they remain rational and of type $(p,p)$ after conjugation. Deligne called a Hodge class \emph{absolute} if it
shares this property. Absolute Hodge classes sit between algebraic classes and Hodge classes, and, unlike algebraic classes, one can prove things about them: Deligne's
theorem that all Hodge classes on abelian varieties are absolute (Theorem~\ref{absolute-hodge:theorem-deligne}) is the strongest general result in the direction of the Hodge conjecture.
References: \cite{deligne-abelian}, \cite{charles-schnell}, \cite{voisin-2007}, and the source \emph{The Hodge Conjecture} (Chapter 8).

\noindent\textbf{Prerequisites.} Chapter~\ref{hodge-conjecture} (The Hodge Conjecture), Chapter~\ref{etale-cohomology} (\'Etale Cohomology), Chapter~\ref{gaga} (Analytic Spaces and GAGA),
Section~\ref{hodge-structures:section-mumford-tate} (Mumford--Tate Groups), Chapter~\ref{known-cases} (Known Cases and Evidence) and Chapter~\ref{mixed-hodge-structures} (Mixed Hodge Structures).

\noindent\textbf{Conventions.} In this chapter $X$ is a nonsingular projective variety over $\CC$ of dimension $n$, and $\ell$ runs over all primes. We write
$\QQ(p) = (2\pi i)^p\QQ \subseteq \CC$ and $H^k_B(X)(p) = H^k(X^{\mathrm{an}}; \QQ) \otimes \QQ(p) \subseteq H^k(X^{\mathrm{an}}; \CC)$ for twisted Betti cohomology, $H^k_{\mathrm{dR}}(X) = H^k_{\mathrm{dR}}(X/\CC)$
(Definition~\ref{gaga:definition-algebraic-de-rham}), and $H^k_\ell(X)(p) = H^k_{\mathrm{et}}(X, \QQ_\ell(p))$ (Definition~\ref{etale-cohomology:definition-l-adic}; here $X$ is already over an algebraically
closed field).

\section{Conjugate varieties and comparison isomorphisms}\label{absolute-hodge:section-conjugates}

\begin{definition}\label{absolute-hodge:definition-conjugate}
Let $\sigma \in \mathrm{Aut}(\CC)$. The \emph{conjugate variety} is $X^\sigma = X \times_{\Spec\CC,\sigma}\Spec\CC$, the base change of $X$ along the morphism $\Spec\CC \to \Spec\CC$ induced by $\sigma$. The
projection $\sigma_X \colon X^\sigma \to X$ is an isomorphism of schemes, but not of $\CC$-schemes. If $X \subseteq \PP^N$ is cut out by homogeneous polynomials $f_j$, then $X^\sigma$ is cut out by the
polynomials $f_j^\sigma$ obtained by applying $\sigma$ to the coefficients, and on closed points $\sigma_X^{-1}$ is $(x_0 : \cdots : x_N) \mapsto (\sigma(x_0) : \cdots : \sigma(x_N))$.
\end{definition}

\begin{example}\label{absolute-hodge:example-elliptic}
Let $E$ be the elliptic curve $y^2 = x^3 + ax + b$. Then $E^\sigma$ is $y^2 = x^3 + \sigma(a)x + \sigma(b)$, so $j(E^\sigma) = \sigma(j(E))$, and $E^\sigma \cong E$ if and only if $\sigma$ fixes $j(E)$. If
$j(E)$ is transcendental, then $\sigma(j(E))$ can be any transcendental number, and $E$ has uncountably many pairwise non-isomorphic conjugates. If $E$ is defined over $\QQ$, then $E^\sigma = E$ for
every $\sigma$, but $\sigma_E$ is a discontinuous bijection of $E(\CC)$ for $\sigma \neq \mathrm{id}$, complex conjugation.
\end{example}

\begin{definition}\label{absolute-hodge:definition-transport}
For $\sigma \in \mathrm{Aut}(\CC)$ define
\begin{enumerate}
\item $\sigma_* \colon H^k_{\mathrm{dR}}(X) \to H^k_{\mathrm{dR}}(X^\sigma)$, the $\sigma$-linear isomorphism induced by $\sigma_X^*\Omega^\bullet_{X/\CC} = \Omega^\bullet_{X^\sigma/\CC}$
(Corollary~\ref{gaga:corollary-conjugates}); it maps $F^pH^k_{\mathrm{dR}}(X)$ onto $F^pH^k_{\mathrm{dR}}(X^\sigma)$;
\item $\sigma_* = \sigma_X^* \colon H^k_\ell(X)(p) \to H^k_\ell(X^\sigma)(p)$, the pullback in \'etale cohomology along the isomorphism of schemes $\sigma_X$, which identifies the sheaves $\mu_{\ell^m}$; it is
$\QQ_\ell$-linear.
\end{enumerate}
\end{definition}

\begin{remark}\label{absolute-hodge:remark-no-betti}
There is no map $H^k_B(X) \to H^k_B(X^\sigma)$ for general $\sigma$: $\sigma_X$ is not continuous for the classical topology unless $\sigma$ is the identity or complex conjugation, and
$X^{\mathrm{an}}$ and $(X^\sigma)^{\mathrm{an}}$ need not even be homeomorphic (Remark~\ref{gaga:remark-conjugates-not-homeomorphic}). The \emph{comparison isomorphisms}
\[
c_{\mathrm{dR}} \colon H^k_B(X)(p) \otimes_\QQ\CC \xrightarrow{\ \sim\ } H^k_{\mathrm{dR}}(X), \qquad c_\ell \colon H^k_B(X)(p) \otimes_\QQ\QQ_\ell \xrightarrow{\ \sim\ } H^k_\ell(X)(p)
\]
of Theorem~\ref{gaga:theorem-grothendieck-comparison} and Theorem~\ref{etale-cohomology:theorem-comparison} (the second uses $\ZZ(1) \otimes \ZZ/\ell^m \cong \mu_{\ell^m}$, $2\pi i/\ell^m \mapsto e^{2\pi i/\ell^m}$) give two
further structures on $H^k_{\mathrm{dR}}(X)$ and $H^k_\ell(X)$: a $\QQ$-structure and the complex conjugation it defines. The map $\sigma_*$ preserves the algebraic structure (the Hodge filtration,
cup products, algebraic cycle classes) but has no reason to preserve these transcendental ones. In particular it preserves $F^p$ but not $\overline{F^q}$, so it need not map $H^{p,q}(X)$ to $H^{p,q}(X^\sigma)$.
\end{remark}

\begin{proposition}\label{absolute-hodge:proposition-invariants}
The Betti numbers, the Hodge numbers, the dimension of $\mathrm{Hdg}^1(X)$ and the Picard number of $X$ and $X^\sigma$ agree.
\end{proposition}

\begin{proof}
Betti and Hodge numbers: Corollary~\ref{gaga:corollary-conjugates}. For divisors, $\sigma_X^*$ gives a bijection between divisors on $X$ and divisors on $X^\sigma$ that preserves linear and algebraic
equivalence, since these are defined algebraically; so the N\'eron--Severi groups are isomorphic and the Picard numbers agree. By the Lefschetz $(1,1)$-theorem
(Theorem~\ref{hodge-conjecture:theorem-lefschetz-11}), $\dim\mathrm{Hdg}^1(X)$ is the Picard number.
\end{proof}

\section{Absolute Hodge classes}\label{absolute-hodge:section-absolute}

\begin{definition}[Deligne]\label{absolute-hodge:definition-hodge-cycle}
Write $H^{2p}_{\mathbb{A}}(X)(p) = H^{2p}_{\mathrm{dR}}(X) \times \prod_\ell H^{2p}_\ell(X)(p)$ and $t = (t_{\mathrm{dR}}, (t_\ell)_\ell)$ for its elements.
\begin{enumerate}
\item $t$ is a \emph{Hodge cycle} on $X$ if there is $t_B \in H^{2p}_B(X)(p)$ with $c_{\mathrm{dR}}(t_B) = t_{\mathrm{dR}}$ and $c_\ell(t_B) = t_\ell$ for all $\ell$, and $t_{\mathrm{dR}} \in F^pH^{2p}_{\mathrm{dR}}(X)$.
\item $t$ is an \emph{absolute Hodge cycle} if $\sigma_*t = (\sigma_*t_{\mathrm{dR}}, (\sigma_*t_\ell)_\ell)$ is a Hodge cycle on $X^\sigma$ for every $\sigma \in \mathrm{Aut}(\CC)$.
\end{enumerate}
The absolute Hodge cycles form a $\QQ$-vector space $C^p_{\mathrm{AH}}(X)$: $\sigma_*$ is additive and fixes rational scalars.
\end{definition}

\begin{lemma}\label{absolute-hodge:lemma-hodge-cycles}
The map $t \mapsto (2\pi i)^{-p}t_B$ is an isomorphism from the $\QQ$-vector space of Hodge cycles on $X$ onto $\mathrm{Hdg}^p(X)$. Its inverse sends $\alpha$ to
$(c_{\mathrm{dR}}((2\pi i)^p\alpha), (c_\ell((2\pi i)^p\alpha))_\ell)$.
\end{lemma}

\begin{proof}
$t_B$ is determined by $t$, because $c_\ell$ is injective. Put $\alpha = (2\pi i)^{-p}t_B \in H^{2p}(X^{\mathrm{an}}; \QQ)$. Since $F^p$ is a complex subspace, $t_{\mathrm{dR}} \in F^p$ if and only if $\alpha \in F^p$.
A rational class is fixed by complex conjugation, so $\alpha \in F^p$ if and only if $\alpha \in F^p \cap \overline{F^p} = H^{p,p}$ (Proposition~\ref{hodge-structures:proposition-filtration}), that is,
$\alpha \in \mathrm{Hdg}^p(X)$. The inverse map is clear.
\end{proof}

We write $\mathrm{Hdg}^p_{\mathrm{AH}}(X) \subseteq \mathrm{Hdg}^p(X)$ for the image of $C^p_{\mathrm{AH}}(X)$, and call its elements \emph{absolute Hodge classes}.

\begin{proposition}\label{absolute-hodge:proposition-algebraic}
Classes of algebraic cycles are absolute Hodge: $Z^p_{\mathrm{alg}}(X) \subseteq \mathrm{Hdg}^p_{\mathrm{AH}}(X) \subseteq \mathrm{Hdg}^p(X)$.
\end{proposition}

\begin{proof}
Let $Z$ be a cycle of codimension $p$ with $\QQ$-coefficients. There are cycle classes $\mathrm{cl}_{\mathrm{dR}}(Z) \in F^pH^{2p}_{\mathrm{dR}}(X)$ and $\mathrm{cl}_\ell(Z) \in H^{2p}_\ell(X)(p)$
(Theorem~\ref{etale-cohomology:theorem-l-adic-properties}(4)), both defined algebraically, and with suitable sign conventions they are the images of $(2\pi i)^p\mathrm{cl}(Z)$ under $c_{\mathrm{dR}}$ and $c_\ell$:
for $\ell$-adic cohomology this is Theorem~\ref{etale-cohomology:theorem-comparison}, and for de Rham cohomology it is \cite{deligne-abelian}. So $\mathrm{cl}_{\mathbb{A}}(Z) = (\mathrm{cl}_{\mathrm{dR}}(Z), (\mathrm{cl}_\ell(Z))_\ell)$ is a Hodge cycle. Since the cycle classes are defined algebraically,
they are compatible with transport of structure along the isomorphism of schemes $\sigma_X$: $\sigma_*\mathrm{cl}_{\mathbb{A}}(Z) = \mathrm{cl}_{\mathbb{A}}(Z^\sigma)$, where $Z^\sigma = \sigma_X^{-1}(Z)$ is a cycle on
$X^\sigma$. By the same argument on $X^\sigma$, this is a Hodge cycle. The second inclusion is the case $\sigma = \mathrm{id}$.
\end{proof}

\begin{remark}\label{absolute-hodge:remark-two-pi-i}
The factor $(2\pi i)^p$ is forced: already for a point on $\PP^1$, the de Rham class is $\pm 2\pi i$ times the Betti class (Exercise~\ref{absolute-hodge:exercise-two-pi-i}). This is why the rational
structure relevant for cycle classes of codimension $p$ is $H^{2p}_B(X)(p)$ and not $H^{2p}(X^{\mathrm{an}}; \QQ)$.
\end{remark}

\begin{conjecture}[Absolute Hodge conjecture]\label{absolute-hodge:conjecture-absolute}
Every Hodge class is absolute Hodge: $\mathrm{Hdg}^p_{\mathrm{AH}}(X) = \mathrm{Hdg}^p(X)$.
\end{conjecture}

By Proposition~\ref{absolute-hodge:proposition-algebraic}, the Hodge conjecture for $X$ and $X^\sigma$ in all codimensions implies Conjecture~\ref{absolute-hodge:conjecture-absolute}. The
converse is not expected to be provable directly: being absolute Hodge is a condition on cohomology and its comparison isomorphisms, and it never produces a subvariety.

\begin{remark}\label{absolute-hodge:remark-de-rham-betti}
Several variants are in use. Omitting the $\ell$-adic components gives \emph{de Rham--Betti absolute} classes, which is the version in the source \emph{The Hodge Conjecture}; Deligne's theorem below holds
for the full definition, and all known results hold for both. Andr\'e's \emph{motivated cycles} \cite{andre-1996-motifs} are a smaller class, generated by algebraic cycles and the inverse Lefschetz
operators (Corollary~\ref{absolute-hodge:corollary-lambda}); they are absolute Hodge.
\end{remark}

\section{Operations on absolute Hodge classes}\label{absolute-hodge:section-operations}

\begin{proposition}\label{absolute-hodge:proposition-operations}
Let $f \colon X \to Y$ be a morphism of nonsingular projective varieties, of relative dimension $d = \dim X - \dim Y$. Then absolute Hodge classes are preserved by
\begin{enumerate}
\item pullback $f^* \colon \mathrm{Hdg}^p(Y) \to \mathrm{Hdg}^p(X)$ and pushforward $f_* \colon \mathrm{Hdg}^p(X) \to \mathrm{Hdg}^{p-d}(Y)$;
\item cup product, so $\bigoplus_p\mathrm{Hdg}^p_{\mathrm{AH}}(X)$ is a $\QQ$-subalgebra of $\bigoplus_p\mathrm{Hdg}^p(X)$;
\item the K\"unneth projections $H^{2p}(X \times Y) \to H^i(X) \otimes H^{2p-i}(Y)$: the K\"unneth components of an absolute Hodge class on $X \times Y$ are absolute Hodge;
\item correspondences: if $\gamma \in \mathrm{Hdg}^q_{\mathrm{AH}}(X \times Y)$ then $\gamma_* = \mathrm{pr}_{2*}(\mathrm{pr}_1^*(-) \cup \gamma)$ maps absolute Hodge classes on $X$ to absolute Hodge classes on $Y$.
\end{enumerate}
\end{proposition}

\begin{proof}
Each operation is defined in Betti, de Rham and $\ell$-adic cohomology. In de Rham and $\ell$-adic cohomology it is defined algebraically, from pullbacks, cup products and trace maps, so it commutes with
$\sigma_*$: for instance $\sigma_*(f^*t) = (f^\sigma)^*(\sigma_*t)$, where $f^\sigma \colon X^\sigma \to Y^\sigma$ is the conjugate morphism. The operations are compatible with the comparison isomorphisms
$c_{\mathrm{dR}}$ and $c_\ell$ (with the Tate twists indicated: $f_*$ goes from $H^{2p}(X)(p)$ to $H^{2p-2d}(Y)(p - d)$), and on Betti cohomology they are morphisms of Hodge structures, so they map Hodge cycles to Hodge
cycles \cite{deligne-abelian}. Given an absolute Hodge cycle $t$ on $Y$ and $\sigma \in \mathrm{Aut}(\CC)$, the cycle $\sigma_*(f^*t) = (f^\sigma)^*(\sigma_*t)$ is the image of the Hodge cycle $\sigma_*t$ on $Y^\sigma$
under an operation that preserves Hodge cycles, hence a Hodge cycle. The other cases are identical: for (3), the K\"unneth decomposition is built from $\mathrm{pr}_1^*$, $\mathrm{pr}_2^*$ and cup product, so the
projections commute with $\sigma_*$ and preserve Hodge cycles, because the K\"unneth decomposition of $H^{2p}(X \times Y)$ is a decomposition into sub-Hodge structures; and (4) is a composite of (1) and (2).
\end{proof}

\begin{corollary}\label{absolute-hodge:corollary-kunneth}
The K\"unneth components $\pi_i \in H^{2n-i}(X) \otimes H^i(X) \subseteq H^{2n}(X \times X)$ of the class of the diagonal are absolute Hodge classes.
\end{corollary}

\begin{proof}
The class of the diagonal is algebraic, hence absolute Hodge by Proposition~\ref{absolute-hodge:proposition-algebraic}; apply Proposition~\ref{absolute-hodge:proposition-operations}(3).
\end{proof}

\begin{corollary}\label{absolute-hodge:corollary-lambda}
Let $h \in H^2(X)$ be the class of an ample divisor and $\Lambda$ the inverse Lefschetz operator of $L = h \cup -$ (Proposition~\ref{kahler:proposition-lambda-unique}). The class
$\lambda \in H^{2n-2}(X \times X)(n-1)$ of $\Lambda$, obtained through the K\"unneth decomposition and Poincar\'e duality, is an absolute Hodge class.
\end{corollary}

\begin{proof}
Let $\sigma \in \mathrm{Aut}(\CC)$ and let $h^\sigma$ be the class of the conjugate divisor on $X^\sigma$, with operators $L^\sigma$ and $\Lambda^\sigma$. In de Rham cohomology, $\sigma_*$ intertwines $L$ with $L^\sigma$ (it
preserves cup products and cycle classes) and preserves degrees. Hence $\sigma_*\Lambda\sigma_*^{-1}$ satisfies the equations $[L^\sigma, f] = h_{\deg}$, $[h_{\deg}, f] = -2f$, where $h_{\deg}$ acts by $k - n$ in degree $k$,
and by the uniqueness in Proposition~\ref{kahler:proposition-lambda-unique} it equals $\Lambda^\sigma_{\mathrm{dR}}$. The same argument applies in $\ell$-adic cohomology, using hard Lefschetz there
(Theorem~\ref{etale-cohomology:theorem-l-adic-properties}(6)). On the other hand, the Betti operator $\Lambda^\sigma_B$ is defined over $\QQ$ (Proposition~\ref{kahler:proposition-lambda-unique}),
and its images under $c_{\mathrm{dR}}$ and $c_\ell$ satisfy the defining equations, so they are $\Lambda^\sigma_{\mathrm{dR}}$ and $\Lambda^\sigma_\ell$ by uniqueness again. Finally, $\Lambda^\sigma_B$ is a morphism of Hodge
structures of type $(-1,-1)$. Indeed, let $h(z)$, $z \in \mathbb{S}(\RR) = \CC^*$, be the action of the Deligne torus on $H^\bullet(X^\sigma; \RR)$ (Proposition~\ref{hodge-structures:proposition-deligne-torus}). Since
$L^\sigma$ maps $H^{p,q}$ to $H^{p+1,q+1}$, $h(z)L^\sigma h(z)^{-1} = |z|^{-2}L^\sigma$, and $h(z)$ commutes with the grading operator. So $|z|^{-2}h(z)\Lambda^\sigma_Bh(z)^{-1}$ solves the defining equations for
$L^\sigma$, and by uniqueness $h(z)\Lambda^\sigma_Bh(z)^{-1} = |z|^2\Lambda^\sigma_B$ for all $z$, which says exactly that $\Lambda^\sigma_B$ maps $H^{p,q}$ to $H^{p-1,q-1}$. So the class of $\Lambda^\sigma$ is a Hodge class on $X^\sigma \times X^\sigma$
(Proposition~\ref{hodge-structures:proposition-hodge-classes-morphisms}), and $\sigma_*\lambda$ is a Hodge cycle.
\end{proof}

\begin{remark}\label{absolute-hodge:remark-lefschetz-standard}
Grothendieck's Lefschetz standard conjecture asks for $\lambda$ to be algebraic (Section~\ref{absolute-hodge:section-standard-conjectures}). Corollary~\ref{absolute-hodge:corollary-lambda} proves the weaker statement that $\lambda$ is absolute
Hodge, and the Hodge conjecture would bridge the gap. This is typical: statements whose algebraic form is open are often theorems in the absolute Hodge setting.
\end{remark}

\section{Galois action and fields of definition}\label{absolute-hodge:section-galois}

Every $X$ is defined over a subfield of $\CC$ finitely generated over $\QQ$, since it is cut out by finitely many polynomials: $X = X_0 \times_k\CC$ with $k \subseteq \CC$ finitely generated over $\QQ$ and $X_0$ a
nonsingular projective $k$-variety. Let $\bar k$ be the algebraic closure of $k$ in $\CC$ and $G_k = \mathrm{Gal}(\bar k/k)$. For $\tau \in \mathrm{Aut}(\CC/k)$ we have $X^\tau = X$ canonically, so $\tau_*$ acts on
$H^{2p}_{\mathbb{A}}(X)(p)$. On $H^{2p}_\ell(X)(p) \cong H^{2p}_{\mathrm{et}}(X_0 \times_k\bar k, \QQ_\ell(p))$, where the isomorphism is invariance of \'etale cohomology of proper varieties under extension of
algebraically closed fields \cite[Chapter VI]{milne-etale}, the action of $\tau$ is the Galois action of $\tau|_{\bar k} \in G_k$.

\begin{proposition}\label{absolute-hodge:proposition-galois-finite}
In this situation, $\mathrm{Aut}(\CC/k)$ maps $C^p_{\mathrm{AH}}(X)$ to itself, and it acts on it through a finite quotient of $G_k$. Consequently there is a finite extension $k' \subseteq \bar k$ of $k$ such that
for every absolute Hodge cycle $t$, the $\ell$-adic components $t_\ell \in H^{2p}_{\mathrm{et}}(X_0 \times_k\bar k, \QQ_\ell(p))$ are fixed by $G_{k'}$.
\end{proposition}

\begin{proof}
If $t$ is absolute Hodge and $\tau \in \mathrm{Aut}(\CC/k)$, then $\tau_*t$ is a Hodge cycle on $X^\tau = X$, and $\sigma_*(\tau_*t) = (\sigma\tau)_*t$ is a Hodge cycle on $X^{\sigma\tau}$ for every $\sigma$; so $\tau_*t$ is
absolute Hodge. The space $V = C^p_{\mathrm{AH}}(X)$ is finite-dimensional over $\QQ$, since $t \mapsto t_B$ embeds it in $H^{2p}_B(X)(p)$. Fix a prime $\ell$. The map $V \otimes \QQ_\ell \to H^{2p}_\ell(X)(p)$,
$t \otimes a \mapsto at_\ell$, is injective, because it is the restriction of the comparison isomorphism $c_\ell$ to $V_B \otimes \QQ_\ell$, where $V_B = \{t_B : t \in V\}$. It is equivariant, and the action of
$\tau$ on the right depends only on $\tau|_{\bar k}$. So the action on $V$ is given by a homomorphism $\rho \colon G_k \to \mathrm{GL}(V)$, and $\rho$ is continuous for the $\ell$-adic topology on
$\mathrm{GL}(V \otimes \QQ_\ell)$, because the Galois action on $\ell$-adic cohomology is continuous. The image $\rho(G_k)$ is therefore a compact subgroup of $\mathrm{GL}(V \otimes \QQ_\ell)$, and it is contained in
the countable group $\mathrm{GL}(V)$. A countable compact Hausdorff group is finite: by the Baire category theorem one of its points is not nowhere dense, hence is open, so by homogeneity every
point is open, and a discrete compact space is finite. Hence $\ker\rho$ is an open subgroup of $G_k$, of the form $G_{k'}$ with $k'/k$ finite, and it fixes every $t$, hence every $t_\ell$.
\end{proof}

\begin{corollary}\label{absolute-hodge:corollary-tate-classes}
If $X = X_0 \times_k\CC$ as above, the $\ell$-adic components of absolute Hodge classes are \emph{Tate classes}: classes in $H^{2p}_{\mathrm{et}}(X_0 \times_k\bar k, \QQ_\ell(p))$ fixed by an open subgroup
of $G_k$. If the absolute Hodge conjecture holds for $X$, then every Hodge class of $X$ gives, through $c_\ell$, a Tate class.
\end{corollary}

\begin{proof}
The first statement is Proposition~\ref{absolute-hodge:proposition-galois-finite}. For the second, a Hodge class $\alpha$ determines the Hodge cycle of Lemma~\ref{absolute-hodge:lemma-hodge-cycles}, which is
absolute Hodge by hypothesis, and its $\ell$-adic component is $c_\ell((2\pi i)^p\alpha)$.
\end{proof}

\begin{remark}\label{absolute-hodge:remark-hodge-to-tate}
Corollary~\ref{absolute-hodge:corollary-tate-classes} is the bridge from Hodge theory to arithmetic. Algebraic classes of cycles defined over $k'$ are Tate classes for $G_{k'}$
(Theorem~\ref{etale-cohomology:theorem-l-adic-properties}(4)), and the Tate conjecture predicts the converse (Section~\ref{absolute-hodge:section-tate-conjecture}). Absolute Hodge classes give Tate classes without
producing cycles.
\end{remark}

\section{Principle B}\label{absolute-hodge:section-principle-b}

\begin{theorem}[Deligne's Principle B]\label{absolute-hodge:theorem-principle-b}
Let $f \colon \mathcal{X} \to S$ be a smooth projective morphism of nonsingular complex varieties with $S$ connected, and let $t$ be a global section of the local system $R^{2p}f_*\QQ(p)$ such that
$t_s \in H^{2p}_B(X_s)(p)$ is (the Betti component of) a Hodge cycle for every $s \in S$. If $t_{s_0}$ is absolute Hodge for one point $s_0 \in S$, then $t_s$ is absolute Hodge for every $s \in S$.
\end{theorem}

\begin{proof}[Sketch of proof]
Fix $\sigma \in \mathrm{Aut}(\CC)$. The relative de Rham cohomology $\mathcal{H}^{2p}_{\mathrm{dR}}(\mathcal{X}/S)$ is an algebraic vector bundle on $S$ with an algebraic flat connection, the Gauss--Manin connection
(Definition~\ref{vhs:definition-gauss-manin}), which is defined algebraically by Katz and Oda \cite{katz-oda-1968}; hence conjugation by $\sigma$ carries it to the Gauss--Manin connection of
$f^\sigma \colon \mathcal{X}^\sigma \to S^\sigma$. We use three facts. First, the de Rham components $t_{\mathrm{dR},s}$ form a flat section of $\mathcal{H}^{2p}_{\mathrm{dR}}$, and flat sections of the Gauss--Manin
connection are algebraic, by its regularity; so $\sigma_*t_{\mathrm{dR}}$ is a flat algebraic section over $S^\sigma$. Second, the $\ell$-adic components form a global section of the \'etale local system
$R^{2p}f_*\QQ_\ell(p)$, and $\sigma_*$ carries it to a global section over $S^\sigma$. We omit the proofs of these two facts \cite{deligne-abelian}. Third, by hypothesis, $u = \sigma_*t$ is a Hodge cycle at the point
$\sigma(s_0) \in S^\sigma$, with Betti component $u_B$.

Since $u_\ell$ is a global section, it is invariant under the monodromy of $S^\sigma$; as $c_\ell$ is compatible with monodromy, so is $u_B \otimes 1$, and hence $u_B$ is the value at $\sigma(s_0)$ of a global section
of $R^{2p}f^\sigma_*\QQ(p)$. Its de Rham and $\ell$-adic incarnations are flat and agree with $u_{\mathrm{dR}}$ and $u_\ell$ at $\sigma(s_0)$, hence everywhere, because flat sections of a local system on a
connected space are determined by one value. So $u$ is rational at every point of $S^\sigma$. By the global invariant cycle theorem (Theorem~\ref{mixed-hodge-structures:theorem-global-invariant}), the global sections
of $R^{2p}f^\sigma_*\QQ(p)$ form a Hodge structure $V$ that maps to each fibre by an injective morphism of Hodge structures: $V$ is the quotient of $H^{2p}$ of a nonsingular projective compactification of
$\mathcal{X}^\sigma$ by the kernel of the restriction map to global sections, which does not depend on the fibre. The element $u_B \in V$ is of type $(0,0)$ in the fibre at $\sigma(s_0)$, hence in $V$, hence in
every fibre. Thus $u = \sigma_*t$ is a Hodge cycle at every point of $S^\sigma$. As $\sigma$ was arbitrary, $t_s$ is absolute Hodge for every $s$.
\end{proof}

\begin{remark}\label{absolute-hodge:remark-variational}
Principle B is the absolute Hodge analogue of the variational Hodge conjecture (Conjecture~\ref{known-cases:conjecture-variational}), which is open for algebraic classes. The difference is that absolute Hodge
cycles are defined by conditions on cohomology, which propagate along flat sections, while algebraic cycles are geometric objects whose existence on one fibre does not obviously imply existence on the
others.
\end{remark}

\section{Deligne's theorem for abelian varieties}\label{absolute-hodge:section-deligne}

\begin{theorem}[Deligne]\label{absolute-hodge:theorem-deligne}
Let $A$ be a complex abelian variety. Every Hodge class on $A$ is absolute Hodge: $\mathrm{Hdg}^p_{\mathrm{AH}}(A) = \mathrm{Hdg}^p(A)$ for all $p$. The same holds for Hodge classes in all tensor spaces
$T^{m,m'}(H^1(A; \QQ))(r)$ of weight $0$, defined analogously.
\end{theorem}

\begin{proof}[Outline of proof]
The proof is \cite{deligne-abelian}; see also \cite{charles-schnell}. It combines four ingredients.
\begin{enumerate}
\item \emph{Principle B} (Theorem~\ref{absolute-hodge:theorem-principle-b}).
\item \emph{Principle A.} If $t_1, \ldots, t_r$ are absolute Hodge tensors on $A$, and $G \subseteq \mathrm{GL}(H^1(A; \QQ))$ is the subgroup fixing them, then every tensor fixed by $G$ is absolute Hodge. Applied
when $G$ is the Hodge group, this reduces the theorem to finding absolute Hodge tensors whose stabiliser is $\mathrm{Hg}(A)$ (Theorem~\ref{hodge-structures:theorem-hodge-classes-invariants}).
\item \emph{Weil classes are absolute Hodge.} Let $K$ be an imaginary quadratic field and consider abelian varieties of Weil type for $K$ (Proposition~\ref{known-cases:proposition-weil-classes} treats
dimension $4$; the construction is the same in every even dimension), with a fixed type of polarisation. They form a connected family over which the Weil classes form a flat family of Hodge classes. For
the \emph{split} type, the family contains abelian varieties of the form $B \otimes_\QQ K$, on which the Weil classes are algebraic. By Principle B the Weil classes are absolute Hodge on every member.
\item \emph{Reduction to CM type, and the CM case.} A Hodge class $t$ on $A$ is fixed by $\mathrm{Hg}(A)$, hence remains Hodge on every abelian variety in the connected family parametrised by the
Mumford--Tate domain of $A$, and this family contains abelian varieties of CM type. By Principle B it suffices to treat $A$ of CM type. For those, the Hodge group is a torus
(Proposition~\ref{hodge-structures:proposition-cm}), and one shows that it is the stabiliser of the pullbacks of split Weil classes along homomorphisms from $A$ to abelian varieties of Weil type, together with
divisor classes. Principle A and (3) conclude.
\end{enumerate}
\end{proof}

\begin{remark}\label{absolute-hodge:remark-force-and-limits}
Theorem~\ref{absolute-hodge:theorem-deligne} says that a Hodge class on an abelian variety, however transcendentally it was produced, behaves under all automorphisms of $\CC$ as if it were algebraic. Its
limit is equally clear: the Weil classes on a general abelian fourfold of Weil type are absolute Hodge but, apart from the cases listed in Chapter~\ref{known-cases}, not known to be algebraic. Deligne used
the theorem, with Milne, to build an unconditional Tannakian category of motives for absolute Hodge cycles containing the motives of abelian varieties \cite{deligne-milne}. A deformation argument of the
same kind, combined with the Kuga--Satake construction, had earlier given Deligne's proof of the Weil conjectures for K3 surfaces \cite{deligne-k3}.
\end{remark}

\section{Fields of definition of Hodge loci}\label{absolute-hodge:section-hodge-loci}

\begin{lemma}\label{absolute-hodge:lemma-countable-conjugates}
Let $Z \subseteq \PP^N_\CC$ be a closed subvariety. If the set $\{\sigma(Z) : \sigma \in \mathrm{Aut}(\CC/\bar{\QQ})\}$ is countable, then $Z$ is defined over $\bar{\QQ}$.
\end{lemma}

\begin{proof}
Fix a monomial order. The homogeneous ideal of $Z$ has a unique reduced Gr\"obner basis, and the field $k_0 \subseteq \CC$ generated by its coefficients is the smallest field of definition of $Z$; for
$\sigma \in \mathrm{Aut}(\CC)$ the reduced Gr\"obner basis of $\sigma(Z)$ is obtained by applying $\sigma$ to the coefficients, so $\sigma(Z) = \tau(Z)$ if and only if $\sigma|_{k_0} = \tau|_{k_0}$. If $k_0 \not\subseteq \bar{\QQ}$,
it contains a transcendental element $c$, and for each transcendental $c' \in \CC$ there is $\sigma \in \mathrm{Aut}(\CC/\bar{\QQ})$ with $\sigma(c) = c'$ (extend $c \mapsto c'$ from $\bar{\QQ}(c)$ to $\CC$, using transcendence
bases and algebraic closures). The values $\sigma(c)$, and hence the restrictions $\sigma|_{k_0}$ and the conjugates $\sigma(Z)$, are then uncountable, a contradiction. So $k_0 \subseteq \bar{\QQ}$.
\end{proof}

\begin{proposition}[Voisin]\label{absolute-hodge:proposition-voisin}
Let $f \colon \mathcal{X} \to S$ be a smooth projective morphism of nonsingular quasi-projective varieties defined over $\bar{\QQ}$. Let $Z \subseteq S$ be an irreducible closed subvariety, $s \in Z$ and
$\alpha \in \mathrm{Hdg}^p(X_s)$ such that the parallel transport of $\alpha$ along every path in $Z$ is a Hodge class, and suppose $Z$ is maximal with this property, that is, an irreducible component of the Hodge
locus (Definition~\ref{vhs:definition-hodge-locus}). If $\alpha$ is absolute Hodge, then $Z$ is defined over $\bar{\QQ}$, and its conjugates are components of Hodge loci.
\end{proposition}

\begin{proof}[Sketch of proof]
$Z$ is algebraic by the Cattani--Deligne--Kaplan theorem (Theorem~\ref{vhs:theorem-cdk}). By Principle B, applied over the nonsingular locus of $Z$, the class is absolute Hodge at every point of $Z$. Let
$\sigma \in \mathrm{Aut}(\CC/\bar{\QQ})$; then $\mathcal{X}^\sigma = \mathcal{X}$ and $S^\sigma = S$. As in the proof of Theorem~\ref{absolute-hodge:theorem-principle-b}, $\sigma$ transports the flat family of absolute
Hodge cycles over $Z$ to a flat family of Hodge cycles over $\sigma(Z)$. So $\sigma(Z)$ is contained in a component $Z'$ of the Hodge locus of a class, and applying the same argument to $\sigma^{-1}$ and $Z'$
shows $\sigma^{-1}(Z') \subseteq Z$, so $Z' = \sigma(Z)$. The Hodge loci of all rational classes have only countably many components in total, since there are countably many rational classes at a base point
and each Hodge locus is an analytic subset with countably many components. So $Z$ has countably many conjugates, and Lemma~\ref{absolute-hodge:lemma-countable-conjugates} (after embedding $S$ in a
projective space over $\bar{\QQ}$ and taking closures) shows that $Z$ is defined over $\bar{\QQ}$ \cite{voisin-2007}.
\end{proof}

\begin{corollary}\label{absolute-hodge:corollary-abelian-loci}
In a family of abelian varieties defined over $\bar{\QQ}$, every component of the Hodge locus is defined over $\bar{\QQ}$.
\end{corollary}

\begin{proof}
Combine Theorem~\ref{absolute-hodge:theorem-deligne} with Proposition~\ref{absolute-hodge:proposition-voisin}.
\end{proof}

\begin{remark}\label{absolute-hodge:remark-hodge-loci}
Voisin also showed a partial converse: if a component of a Hodge locus is defined over $\bar{\QQ}$ and its conjugates are components of Hodge loci, then the Hodge classes over it are absolute Hodge,
under a mild condition on the variation of Hodge structure \cite{voisin-2007}. So the absolute Hodge conjecture is, to a large extent, a statement about the arithmetic of Hodge loci. For families of
abelian varieties Corollary~\ref{absolute-hodge:corollary-abelian-loci} reflects the fact that special subvarieties of Shimura varieties are defined over $\bar{\QQ}$.
\end{remark}

\section{Exercises}\label{absolute-hodge:section-exercises}

\begin{exercise}\label{absolute-hodge:exercise-picard}
Let $X$ be a nonsingular projective surface and $\sigma \in \mathrm{Aut}(\CC)$. Show that $\sigma_*$ maps $\mathrm{Hdg}^1(X)$, viewed through Lemma~\ref{absolute-hodge:lemma-hodge-cycles}, isomorphically onto
$\mathrm{Hdg}^1(X^\sigma)$, and that $\mathrm{Hdg}^2(X) = H^4(X; \QQ)$ also consists of absolute Hodge classes. Conclude that the absolute Hodge conjecture holds for surfaces.
\end{exercise}

\begin{solution}
By the Lefschetz $(1,1)$-theorem, $\mathrm{Hdg}^1(X)$ is spanned by divisor classes, which are absolute Hodge (Proposition~\ref{absolute-hodge:proposition-algebraic}), and $\sigma_*\mathrm{cl}(D) = \mathrm{cl}(D^\sigma)$.
The same applies to $\sigma^{-1}$ on $X^\sigma$, so the map is bijective. $H^4(X; \QQ)$ is spanned by the class of a point, which is algebraic. $\mathrm{Hdg}^0 = H^0$ is spanned by the class of $X$. So all Hodge
classes on $X$ are algebraic, hence absolute Hodge.
\end{solution}

\begin{exercise}\label{absolute-hodge:exercise-two-pi-i}
On $\PP^1$ with coordinate $z$, the line bundle $\mathcal{O}(1)$ has transition function $z$ on $U_0 \cap U_1$, and its de Rham Chern class is represented in \v{C}ech--de Rham cohomology by the closed
form $dz/z$ on $U_0 \cap U_1$. Using a partition of unity, show that the corresponding smooth $2$-form $\eta$ satisfies $\int_{\PP^1}\eta = \pm 2\pi i$. Deduce that the de Rham class of a point is $\pm 2\pi i$ times the
Betti class, so that it lies in $H^2_B(\PP^1)(1)$.
\end{exercise}

\begin{solution}
Let $\rho_0$ be a smooth function on $\PP^1$ equal to $1$ for $|z| \leq 1$ and to $0$ for $|z| \geq 2$, and $\rho_1 = 1 - \rho_0$. The \v{C}ech cocycle $dz/z$ is the difference of the $1$-forms $\rho_0\,dz/z$ on $U_1$ and
$-\rho_1\,dz/z$ on $U_0$ (each extended by zero), and the corresponding global $2$-form is $\eta = d(\rho_0\,dz/z)$, which equals $-d(\rho_1\,dz/z)$ where both are defined, since $d(dz/z) = 0$. It vanishes for $|z| \leq 1$
and for $|z| \geq 2$. By Stokes' theorem on the annulus $\frac12 \leq |z| \leq 3$, $\int\eta = \int_{|z|=3}\rho_0\,dz/z - \int_{|z|=1/2}\rho_0\,dz/z = 0 - 2\pi i$. The sign depends on the conventions for the
\v{C}ech--de Rham isomorphism. As $\int_{\PP^1}$ of the Betti class of a point is $1$, the de Rham class is $\pm 2\pi i$ times the Betti class.
\end{solution}

\begin{exercise}\label{absolute-hodge:exercise-correspondence}
Let $X$ and $Y$ be nonsingular projective varieties and $\gamma \in \mathrm{Hdg}^q(X \times Y)$ an absolute Hodge class such that $\gamma_* \colon H^{2p}(X; \QQ) \to H^{2p+2q-2\dim X}(Y; \QQ)$ is surjective.
Show that if the absolute Hodge conjecture holds for $X$ in codimension $p$, then it holds for $Y$ in codimension $p + q - \dim X$.
\end{exercise}

\begin{solution}
$\gamma_*$ is a surjective morphism of polarisable Hodge structures (up to a Tate twist), so by semisimplicity (Theorem~\ref{hodge-structures:theorem-semisimple}) every Hodge class on $Y$ in that degree is
$\gamma_*$ of a Hodge class on $X$, as in the proof of Proposition~\ref{mixed-hodge-structures:proposition-localisation}. By hypothesis that class is absolute Hodge, and $\gamma_*$ preserves absolute Hodge classes
(Proposition~\ref{absolute-hodge:proposition-operations}(4)).
\end{solution}

\begin{exercise}\label{absolute-hodge:exercise-product-abelian}
Let $C$ be a nonsingular projective curve and $J$ its Jacobian. Using the Abel--Jacobi map $C^g \to J$ and Theorem~\ref{absolute-hodge:theorem-deligne}, show that every Hodge class on $C^m$, for every
$m \geq 1$, is absolute Hodge.
\end{exercise}

\begin{solution}
Fix a base point and let $u \colon C \to J$ be the Abel--Jacobi map; $u^* \colon H^1(J; \QQ) \to H^1(C; \QQ)$ is an isomorphism of Hodge structures. By the K\"unneth formula, $H^\bullet(C^m)$ is the direct sum of the
sub-Hodge structures $H^{a_1}(C) \otimes \cdots \otimes H^{a_m}(C)$, $a_i \in \{0, 1, 2\}$, and a Hodge class on $C^m$ is a sum of Hodge classes in these summands, which are absolute Hodge if the summands are,
by Proposition~\ref{absolute-hodge:proposition-operations}(3). Fix a summand, let $S = \{i : a_i = 1\}$ and $T = \{i : a_i = 2\}$. Since $H^0(C)$ is spanned by $1$ and $H^2(C)$ by the class of a point, a Hodge class in
the summand has the form $\mathrm{pr}_S^*(\gamma) \cup \prod_{i \in T}\mathrm{pr}_i^*[\mathrm{pt}]$ with $\gamma$ a Hodge class in the summand $H^1(C)^{\otimes S}$ of $H^{|S|}(C^S)$. The pullback $(u^S)^*$ maps the
corresponding K\"unneth summand $H^1(J)^{\otimes S}$ of $H^{|S|}(J^S)$ isomorphically onto $H^1(C)^{\otimes S}$, as a morphism of Hodge structures, so $\gamma = (u^S)^*\delta$ with $\delta$ a Hodge class on the abelian
variety $J^S$. By Theorem~\ref{absolute-hodge:theorem-deligne} $\delta$ is absolute Hodge, and so is $\gamma$ by Proposition~\ref{absolute-hodge:proposition-operations}(1). The classes $\mathrm{pr}_i^*[\mathrm{pt}]$ are
algebraic, and cup products and pullbacks preserve absolute Hodge classes.
\end{solution}

\section{The Tate Conjecture}\label{absolute-hodge:section-tate-conjecture}
The Tate conjecture is the arithmetic analogue of the Hodge conjecture. The Hodge decomposition is replaced by the action of a Galois group on $\ell$-adic cohomology, and Hodge classes by
\emph{Tate classes}, the classes fixed by an open subgroup of the Galois group. Classes of algebraic cycles defined over the base field are Tate classes, and Tate conjectured the converse for
varieties over finitely generated fields \cite{tate-1965}. Over finite fields there is no Hodge theory and the Tate conjecture is the only statement available; over $\CC$ there is no interesting Galois
group and the Hodge conjecture is the only statement available; over number fields both apply, and the Mumford--Tate conjecture connects them. The best-known evidence is the case of divisors on abelian
varieties, which is Tate's isogeny theorem over finite fields \cite{tate-1966} and Faltings' theorem over number fields \cite{faltings-1983}. References: \cite{tate-1965}, \cite{totaro-2017},
\cite{milne-etale}, and the source \emph{The Hodge Conjecture} (Chapter 8).

\noindent\textbf{Prerequisites.} Chapter~\ref{etale-cohomology} (\'Etale Cohomology), Section~\ref{algebraic-groups:section-abelian-varieties} (Abelian Varieties) and Chapter~\ref{absolute-hodge} (Absolute Hodge Classes).

\noindent\textbf{Conventions.} In this chapter $k$ is a field finitely generated over its prime field, $\bar k$ an algebraic closure, $G_k = \mathrm{Aut}(\bar k/k)$, which is the Galois group of the separable
closure of $k$ in $\bar k$, and $\ell$ a prime different from the characteristic of $k$. $X$ is a nonsingular projective geometrically irreducible variety over $k$, $\bar X = X \times_k\bar k$, and $H^i(\bar X, \QQ_\ell(j))$ is $\ell$-adic cohomology with the continuous
action of $G_k$ (Definition~\ref{etale-cohomology:definition-l-adic}). For a finite extension $k'/k$ inside $\bar k$ we write $G_{k'} \subseteq G_k$ for the corresponding open subgroup.

\section{The statement}\label{absolute-hodge:section-statement}

\begin{definition}\label{absolute-hodge:definition-tate-classes}
The space of \emph{Tate classes of codimension $p$ over $k$} is $T^p_\ell(X/k) = H^{2p}(\bar X, \QQ_\ell(p))^{G_k}$. The space of \emph{geometric Tate classes} is
$T^p_\ell(\bar X) = \bigcup_{k'}T^p_\ell(X_{k'}/k')$, the union over finite extensions $k'/k$, that is, the classes fixed by some open subgroup of $G_k$.
\end{definition}

\begin{lemma}\label{absolute-hodge:lemma-cycles-tate}
Let $Z$ be an algebraic cycle of codimension $p$ on $\bar X$ with $\QQ$-coefficients. Then $Z$ is defined over a finite extension $k'$ of $k$, and $\mathrm{cl}_\ell(Z) \in T^p_\ell(X_{k'}/k')$. If $Z$ is defined
over $k$, then $\mathrm{cl}_\ell(Z) \in T^p_\ell(X/k)$.
\end{lemma}

\begin{proof}
Each component of $Z$ is cut out by finitely many polynomials with coefficients in $\bar k$, which lie in a finite extension $k'$. For $g \in G_{k'}$, $g$ fixes $Z$, and the cycle class is $G_k$-equivariant
(Theorem~\ref{etale-cohomology:theorem-l-adic-properties}(4)), so $g\,\mathrm{cl}_\ell(Z) = \mathrm{cl}_\ell(gZ) = \mathrm{cl}_\ell(Z)$. The same argument with $G_k$ gives the last statement.
\end{proof}

\begin{conjecture}[Tate]\label{absolute-hodge:conjecture-tate}
The cycle class map induces a surjection
\[
\CH^p(X) \otimes_\ZZ\QQ_\ell \longrightarrow H^{2p}(\bar X, \QQ_\ell(p))^{G_k}.
\]
We write $T^p(X/k)$ for this statement, and $T^p(\bar X)$ for the statement that the geometric Tate classes are spanned over $\QQ_\ell$ by classes of cycles on $\bar X$.
\end{conjecture}

\begin{lemma}\label{absolute-hodge:lemma-descent}
If $k'/k$ is a finite extension and $T^p(X_{k'}/k')$ holds, then $T^p(X/k)$ holds. In particular $T^p(\bar X)$ is equivalent to $T^p(X_{k'}/k')$ for all finite $k'/k$.
\end{lemma}

\begin{proof}
Let $k_s$ be the separable closure of $k$ in $k'$, so that $k'/k_s$ is purely inseparable and $G_{k'} = G_{k_s}$. First we show $T^p(X_{k_s}/k_s)$. The projection $\pi \colon X_{k'} \to X_{k_s}$ is finite and a
homeomorphism, so for every subvariety $W \subseteq X_{k'}$ the preimage of $\pi(W)$ is $W$ as a set, and $\pi^*[\pi(W)] = m[W]$ for some integer $m \geq 1$. As $\bar X$ is the same for $k'$ and $k_s$, this gives
$\mathrm{cl}_\ell(W) = \frac1m\mathrm{cl}_\ell(\pi(W))$ in $H^{2p}(\bar X, \QQ_\ell(p))$. So classes of cycles on $X_{k'}$ are combinations of classes of cycles on $X_{k_s}$, and $T^p(X_{k_s}/k_s)$ holds.

Now let $v \in H^{2p}(\bar X, \QQ_\ell(p))^{G_k}$. It is $G_{k_s}$-invariant, so $v = \sum_ia_i\mathrm{cl}_\ell(Z_i)$ with $a_i \in \QQ_\ell$ and $Z_i$ cycles on $X_{k_s}$. Let $g_1, \ldots, g_d$ be representatives of the cosets
$G_k/G_{k_s}$, where $d = [k_s : k]$. On $G_{k_s}$-invariant vectors the operator $\pi_\Gamma = \frac1d\sum_jg_j$ does not depend on the representatives, it fixes $v$, and
$\pi_\Gamma\mathrm{cl}_\ell(Z_i) = \frac1d\mathrm{cl}_\ell(\sum_jg_jZ_i)$. The cycle $\sum_jg_jZ_i$ on $\bar X$ is the pullback of the pushforward of $Z_i$ along the finite separable morphism $X_{k_s} \to X$, hence
the class of a cycle defined over $k$. So $v = \sum_ia_i\pi_\Gamma\mathrm{cl}_\ell(Z_i)$ is a combination of classes of cycles defined over $k$.

For the last statement, a geometric Tate class is fixed by some $G_{k'}$, so $T^p(X_{k'}/k')$ for all $k'$ implies $T^p(\bar X)$. Conversely, if $T^p(\bar X)$ holds and $v$ is a Tate class over $k'$, then
$v = \sum a_i\mathrm{cl}_\ell(Z_i)$ with cycles $Z_i$ on $\bar X$, all defined over some finite extension $k''$ of $k'$ (Lemma~\ref{absolute-hodge:lemma-cycles-tate}). The argument above, with $k'/k$ replaced
by $k''/k'$, used only that $v$ is invariant under the smaller Galois group and is a combination of classes of cycles defined over the larger field; so it shows that $v$ is a combination of classes of
cycles defined over $k'$.
\end{proof}

\begin{counterexample}\label{absolute-hodge:counterexample-algebraically-closed}
The hypothesis that $k$ is finitely generated cannot be dropped. Let $X$ be a nonsingular projective surface over $\CC$ with $h^{2,0}(X) > 0$, for instance a quartic surface in $\PP^3$. For $k = \CC$ the Galois
group is trivial, so every class in $H^2(X, \QQ_\ell(1))$ is invariant. But the $\QQ_\ell$-span of divisor classes has dimension $\rho(X) \leq h^{1,1} < b_2(X)$, by Artin's comparison theorem
(Theorem~\ref{etale-cohomology:theorem-comparison}) and the Lefschetz $(1,1)$-theorem. So $T^1(X/\CC)$ fails. The same happens over $\bar{\QQ}$, and over $\bar{\mathbb{F}}_p$ for a surface with
$\rho(X) < b_2(X)$.
\end{counterexample}

\begin{remark}\label{absolute-hodge:remark-poles}
Over a finite field $k = \mathbb{F}_q$, Tate formulated a companion conjecture: the rank of the group of codimension-$p$ cycles on $X$ modulo numerical equivalence equals the order of the pole of the zeta
function $Z(X, q^{-s})$ at $s = p$ \cite{tate-1965}. Since the eigenvalues of Frobenius on $H^{2p}(\bar X, \QQ_\ell)$ have absolute value $q^p$ \cite{deligne-weil-1}, that order is the multiplicity of the
eigenvalue $q^p$ of the geometric Frobenius on $H^{2p}$. This refinement is tied to $T^p(X/k)$, to the equality of homological and numerical equivalence (Section~\ref{absolute-hodge:section-standard-conjectures}), and to the
semisimplicity of the Frobenius action; see \cite{totaro-2017}.
\end{remark}

\section{Divisors on abelian varieties}\label{absolute-hodge:section-abelian}

For an abelian variety $A$ over $k$, the \emph{Tate module} is $T_\ell A = \varprojlim_mA[\ell^m](\bar k)$, a free $\ZZ_\ell$-module of rank $2\dim A$ with a continuous action of $G_k$, and
$V_\ell A = T_\ell A \otimes_{\ZZ_\ell}\QQ_\ell$. As in Example~\ref{etale-cohomology:example-tate-module}, $H^1(\bar A, \QQ_\ell) \cong \Hom(V_\ell A, \QQ_\ell)$ and $H^i(\bar A, \QQ_\ell) \cong \Lambda^iH^1(\bar A, \QQ_\ell)$,
compatibly with $G_k$ \cite{milne-etale}.

\begin{theorem}[Tate, Zarhin, Faltings]\label{absolute-hodge:theorem-isogeny}
Let $A$ and $B$ be abelian varieties over $k$, where $k$ is finite or of characteristic $\neq 2$. Then:
\begin{enumerate}
\item the natural map $\Hom_k(A, B) \otimes_\ZZ\ZZ_\ell \to \Hom_{G_k}(T_\ell A, T_\ell B)$ is an isomorphism;
\item the representation of $G_k$ on $V_\ell A$ is semisimple.
\end{enumerate}
\end{theorem}

This is due to Tate for finite fields \cite{tate-1966}, Zarhin for finitely generated fields of characteristic $p > 2$ \cite{zarhin-1975}, and Faltings for finitely generated fields of characteristic $0$,
where it is the heart of his proof of the Mordell conjecture \cite{faltings-1983}. Injectivity in (1), and the fact that the cokernel is torsion-free, hold over any field
\cite{mumford-abelian}; the content is surjectivity. Over a finite field the key finiteness input is that there are only finitely many isomorphism classes of polarised abelian varieties of given dimension and
degree over $k$, and over number fields Faltings proves this finiteness with his height.

\begin{corollary}\label{absolute-hodge:corollary-divisors-abelian}
For an abelian variety $A$ over $k$ (as in Theorem~\ref{absolute-hodge:theorem-isogeny}), $T^1(A/k)$ holds: the Tate classes in $H^2(\bar A, \QQ_\ell(1))$ are $\QQ_\ell$-linear combinations of classes of divisors
defined over $k$.
\end{corollary}

\begin{proof}[Sketch of proof]
We use two facts from the theory of abelian varieties \cite{mumford-abelian}. First, $H^2(\bar A, \QQ_\ell(1)) = \Lambda^2H^1(\bar A, \QQ_\ell)(1)$ is the space of alternating $\QQ_\ell$-bilinear forms
$V_\ell A \times V_\ell A \to \QQ_\ell(1)$, and through the Weil pairing $e_\ell \colon V_\ell A \times V_\ell A^\vee \to \QQ_\ell(1)$ such a form is $E_\phi(x, y) = e_\ell(x, \phi y)$ for a unique
$\phi \in \Hom(V_\ell A, V_\ell A^\vee)$ that is symmetric, $\phi^\vee = \phi$ (with the sign conventions of \cite{mumford-abelian}); a form is $G_k$-invariant if and only if $\phi$ is $G_k$-equivariant, because
$e_\ell$ is. Second, the class $c_1(L) \in H^2(\bar A, \QQ_\ell(1))$ of a line bundle corresponds to $\phi = \lambda_L$, the homomorphism $\lambda_L \colon A \to A^\vee$, $a \mapsto t_a^*L \otimes L^{-1}$, and
$L \mapsto \lambda_L$ identifies $\mathrm{NS}(A) \otimes \QQ$ with the symmetric elements of $\Hom_k(A, A^\vee) \otimes \QQ$.

Now let $E$ be a $G_k$-invariant class. Then $E = E_\phi$ with $\phi$ symmetric and equivariant, so by Theorem~\ref{absolute-hodge:theorem-isogeny}(1), applied to $B = A^\vee$,
$\phi = \sum_ja_j\psi_j$ with $a_j \in \QQ_\ell$ and $\psi_j \in \Hom_k(A, A^\vee)$. Replacing $\psi_j$ by $\frac12(\psi_j + \psi_j^\vee)$ we may take the $\psi_j$ symmetric, since $\phi$ is. Each symmetric $\psi_j$ is
$\lambda_{L_j}$ for some $L_j \in \mathrm{NS}(A) \otimes \QQ$, with $L_j$ defined over $k$. Hence $E = \sum_ja_jc_1(L_j)$.
\end{proof}

\begin{example}\label{absolute-hodge:example-products-curves}
Let $C_1$ and $C_2$ be nonsingular projective curves over $k$ with Jacobians $J_1$, $J_2$. By the K\"unneth formula,
$H^2(\bar C_1 \times \bar C_2) = H^2(\bar C_1) \oplus \big(H^1(\bar C_1) \otimes H^1(\bar C_2)\big) \oplus H^2(\bar C_2)$, and the Tate classes in the middle summand, twisted by $(1)$, are the $G_k$-equivariant maps
$V_\ell J_1 \to V_\ell J_2$. By Theorem~\ref{absolute-hodge:theorem-isogeny} these are combinations of homomorphisms $J_1 \to J_2$, which are induced by divisors on $C_1 \times C_2$ (divisorial
correspondences). So $T^1(C_1 \times C_2/k)$ holds. This was Tate's motivating application \cite{tate-1966}.
\end{example}

\section{K3 surfaces and other known cases}\label{absolute-hodge:section-known}

\begin{theorem}\label{absolute-hodge:theorem-known}
The Tate conjecture $T^1$ holds in the following cases.
\begin{enumerate}
\item Abelian varieties over finite fields and over finitely generated fields of characteristic $\neq 2$ (Corollary~\ref{absolute-hodge:corollary-divisors-abelian}), and more generally varieties whose
$H^2$ is controlled by an abelian variety, such as products of curves (Example~\ref{absolute-hodge:example-products-curves}).
\item K3 surfaces over finitely generated fields of characteristic $0$ \cite{andre-1996}.
\item K3 surfaces over finite fields of characteristic $p \geq 3$: for ordinary K3 surfaces by Nygaard \cite{nygaard-1983}, for those of finite height by Nygaard and Ogus \cite{nygaard-ogus-1985}, and in
general by work of Charles \cite{charles-2013}, Maulik \cite{maulik-2014} and Madapusi Pera \cite{madapusi-pera-2015}.
\end{enumerate}
\end{theorem}

The proofs for K3 surfaces use the Kuga--Satake construction (Remark~\ref{known-cases:remark-kuga-satake}), which embeds $H^2$ of a K3 surface into $H^2$ of an abelian variety, and transport
Theorem~\ref{absolute-hodge:theorem-isogeny} through it. The difficulty is arithmetic: the Kuga--Satake abelian variety must be defined over a finitely generated field, and in positive characteristic one needs
integral models of orthogonal Shimura varieties.

\begin{remark}\label{absolute-hodge:remark-artin-tate}
For a surface $X$ over a finite field, Tate showed that $T^1(X/k)$ is equivalent to the finiteness of the $\ell$-primary part of the Brauer group $\mathrm{Br}(X)$, and, with Artin, conjectured a formula
for the special value of the zeta function in terms of $|\mathrm{Br}(X)|$ and the discriminant of the N\'eron--Severi lattice, an analogue of the Birch and Swinnerton-Dyer conjecture
\cite{tate-bourbaki-1966}. The Tate conjecture for divisors thus has a cohomological content (surjectivity of a cycle map) and a finiteness content (Brauer groups) that are equivalent.
\end{remark}

\section{The Mumford--Tate conjecture}\label{absolute-hodge:section-mumford-tate}

Let $k \subseteq \CC$ be a number field and $X$ nonsingular projective over $k$. For each $i$, Artin's comparison theorem identifies $H^i(\bar X, \QQ_\ell)$ with $H^i(X_\CC; \QQ) \otimes \QQ_\ell$
(Theorem~\ref{etale-cohomology:theorem-comparison}, using invariance of \'etale cohomology under extension of algebraically closed fields \cite{milne-etale}). So $G_k$ acts on $H \otimes \QQ_\ell$, where
$H = H^i(X_\CC; \QQ)$ is a Hodge structure with Mumford--Tate group $\mathrm{MT}(H)$ (Definition~\ref{hodge-structures:definition-mumford-tate}).

\begin{definition}\label{absolute-hodge:definition-galois-group}
Let $V$ be a finite-dimensional $\QQ_\ell$-vector space with a continuous representation $\rho \colon G_k \to \mathrm{GL}(V)$. Its \emph{$\ell$-adic algebraic monodromy group} $G_\ell(V)$ is the Zariski closure of
$\rho(G_k)$ in $\mathrm{GL}(V)$, an algebraic group over $\QQ_\ell$, and $G_\ell^0(V)$ is its identity component. For $V = H \otimes \QQ_\ell$ with $H = H^i(X_\CC; \QQ)$ we write $G_\ell$ and $G_\ell^0$.
\end{definition}

\begin{lemma}\label{absolute-hodge:lemma-identity-component}
In the situation of Definition~\ref{absolute-hodge:definition-galois-group}, $G_\ell^0(V)$ does not change if $k$ is replaced by a finite extension, and the vectors of $V$ fixed by some open subgroup of $G_k$ are
exactly the invariants of $G_\ell^0(V)$.
\end{lemma}

\begin{proof}
If $\Gamma_0 \subseteq \Gamma$ are subgroups of $\mathrm{GL}(V)(\QQ_\ell)$ with $\Gamma = \bigcup_{j=1}^rg_j\Gamma_0$, then the Zariski closures satisfy $\bar\Gamma = \bigcup_jg_j\bar\Gamma_0$, so $\bar\Gamma_0$ is a closed subgroup of finite index
in $\bar\Gamma$; it is then a union of connected components of $\bar\Gamma$ and contains $\bar\Gamma^0$. The subgroup $\rho^{-1}(G_\ell^0(V)(\QQ_\ell))$ has finite index in $G_k$ and is closed, hence open, say $G_{k_0}$; the Zariski
closure of $\rho(G_{k_0})$ lies in $G_\ell^0(V)$ and has finite index in $G_\ell(V)$, so it equals $G_\ell^0(V)$. For any finite $k'/k$, the Zariski closure of $\rho(G_{k'})$ has finite index in $G_\ell(V)$, so its identity
component is $G_\ell^0(V)$. This proves the first statement. For the second, the stabiliser of a vector is Zariski closed, so a vector fixed by $\rho(G_{k'})$ is fixed by its Zariski closure, which contains $G_\ell^0(V)$;
conversely a $G_\ell^0(V)$-invariant vector is fixed by $\rho(G_{k_0})$.
\end{proof}

\begin{conjecture}[Mumford--Tate]\label{absolute-hodge:conjecture-mumford-tate}
$G_\ell^0 = \mathrm{MT}(H) \otimes_\QQ\QQ_\ell$ inside $\mathrm{GL}(H \otimes \QQ_\ell)$.
\end{conjecture}

\begin{theorem}[Deligne]\label{absolute-hodge:theorem-deligne-inclusion}
If $X = A$ is an abelian variety and $H = H^1(A_\CC; \QQ)$, then $G_\ell^0 \subseteq \mathrm{MT}(H) \otimes_\QQ\QQ_\ell$.
\end{theorem}

\begin{proof}[Sketch of proof]
Let $\mathrm{MT}' \subseteq \mathrm{GL}(H) \times \mathbb{G}_m$ be the subgroup fixing every Hodge class in every tensor space $T^{m,m'}(H)(r)$ of weight $0$, where $\mathbb{G}_m$ acts on $\QQ(1)$ through the identity
character. Its image under the first projection is $\mathrm{MT}(H)$; this is a reformulation of Theorem~\ref{hodge-structures:theorem-hodge-classes-invariants} which we do not prove \cite{deligne-abelian}. By
Noetherianity, $\mathrm{MT}'$ is the stabiliser of finitely many Hodge classes $t_1, \ldots, t_N$. The spaces $T^{m,m'}(H)(r)$ are Tate twists of Künneth summands of $H^\bullet(A^M_\CC)$ for suitable $M$ (using a
polarisation to identify $H^\vee$ with $H(1)$), so by Theorem~\ref{absolute-hodge:theorem-deligne} the $t_j$ are absolute Hodge, and by Proposition~\ref{absolute-hodge:proposition-galois-finite} their $\ell$-adic
incarnations are fixed by an open subgroup $G_{k'}$. The group $G_k$ acts on $T^{m,m'}(H \otimes \QQ_\ell)(r)$ through $(\rho_\ell, \chi_\ell) \colon G_k \to \mathrm{GL}(H \otimes \QQ_\ell) \times \mathbb{G}_m(\QQ_\ell)$,
$\chi_\ell$ the cyclotomic character. Hence $(\rho_\ell, \chi_\ell)(G_{k'}) \subseteq \mathrm{MT}'(\QQ_\ell)$, so $\rho_\ell(G_{k'}) \subseteq \mathrm{MT}(H)(\QQ_\ell)$, and taking Zariski closures,
$G_\ell^0 \subseteq \mathrm{MT}(H)_{\QQ_\ell}$.
\end{proof}

\begin{proposition}\label{absolute-hodge:proposition-commutants}
Let $A$ be an abelian variety over a number field $k \subseteq \CC$ such that all endomorphisms of $A_\CC$ are defined over $k$. Then $G_\ell$ and $\mathrm{MT}(H)_{\QQ_\ell}$ have the same commutant in
$\End(H \otimes \QQ_\ell)$, namely $\End(A) \otimes \QQ_\ell$.
\end{proposition}

\begin{proof}
By Theorem~\ref{absolute-hodge:theorem-isogeny}(1) with $B = A$, the commutant of $\rho_\ell(G_k)$, which equals the commutant of its Zariski closure $G_\ell$, is $\End_k(A) \otimes \QQ_\ell$, acting on
$V_\ell A$ and dually on $H \otimes \QQ_\ell$. By Corollary~\ref{hodge-structures:corollary-endomorphisms}, the commutant of $\mathrm{MT}(H)$ in $\End(H)$ is $\End_{\mathrm{HS}}(H) = \End(A_\CC) \otimes \QQ$, and taking
commutants commutes with the flat base change $\QQ \to \QQ_\ell$ (it is the solution space of linear equations). As $\End(A_\CC) = \End_k(A)$ by hypothesis, the two commutants agree.
\end{proof}

\begin{remark}\label{absolute-hodge:remark-mumford-tate-known}
Proposition~\ref{absolute-hodge:proposition-commutants} and Theorem~\ref{absolute-hodge:theorem-deligne-inclusion} show that $G_\ell^0 \subseteq \mathrm{MT}_{\QQ_\ell}$ are reductive groups with the same
commutant ($G_\ell^0$ is reductive because $V_\ell A$ is semisimple, Theorem~\ref{absolute-hodge:theorem-isogeny}(2)). This is consistent with the Mumford--Tate conjecture but does not imply it:
reductive subgroups with the same commutant, such as $\mathrm{Sp}_{2g} \subseteq \mathrm{SL}_{2g}$ acting on $\QQ_\ell^{2g}$, can differ. The conjecture is known for elliptic curves: for $E$ without complex
multiplication, Serre's open image theorem \cite{serre-1972} gives $G_\ell = \mathrm{GL}_2$, which is $\mathrm{MT}(H)_{\QQ_\ell}$ by Proposition~\ref{hodge-structures:proposition-elliptic-curves}. It is known for
abelian varieties of CM type, and for many abelian varieties with $\End(A_\CC) = \ZZ$, depending on the dimension \cite{pink-1998}; it is open in general.
\end{remark}

\begin{proposition}\label{absolute-hodge:proposition-hodge-tate}
Let $X$ be a nonsingular projective variety over a number field $k \subseteq \CC$, and $p \geq 0$. Suppose that the geometric Tate classes in $H^{2p}(\bar X, \QQ_\ell(p))$ are exactly the images under $c_\ell$
of the classes $(2\pi i)^p\alpha \otimes a$, $\alpha \in \mathrm{Hdg}^p(X_\CC)$, $a \in \QQ_\ell$; this is what the Mumford--Tate conjecture for $H^{2p}(X_\CC)$ predicts. Then the Hodge conjecture for $X_\CC$
in codimension $p$ is equivalent to $T^p(\bar X)$.
\end{proposition}

\begin{proof}
We use that every algebraic class on $X_\CC$ is the class of a cycle defined over $\bar k$: a cycle on $X_\CC$ is defined over a finitely generated extension of $\bar k$, so it is a member of a family of cycles
parametrised by a variety over $\bar k$, and specialising to a $\bar k$-point does not change its cohomology class, which is locally constant in the family. So the $\QQ$-space $Z^p_{\mathrm{alg}}$ of algebraic
classes is the same for $X_\CC$ and $\bar X$, and $c_\ell$ identifies $Z^p_{\mathrm{alg}} \otimes \QQ_\ell$ with the $\QQ_\ell$-span of the $\ell$-adic cycle classes, by Theorem~\ref{etale-cohomology:theorem-comparison}.
By hypothesis $c_\ell$ identifies $\mathrm{Hdg}^p(X_\CC) \otimes \QQ_\ell$ with $T^p_\ell(\bar X)$. Hence $T^p(\bar X)$ says $Z^p_{\mathrm{alg}} \otimes \QQ_\ell = \mathrm{Hdg}^p \otimes \QQ_\ell$, which, since
$Z^p_{\mathrm{alg}} \subseteq \mathrm{Hdg}^p$ are $\QQ$-subspaces of the same finite-dimensional space, is equivalent to $Z^p_{\mathrm{alg}} = \mathrm{Hdg}^p$.
\end{proof}

\begin{remark}\label{absolute-hodge:remark-hodge-versus-tate}
Neither the Hodge nor the Tate conjecture is known to imply the other. The Hodge conjecture says nothing over finite fields, and the Tate conjecture over number fields is linked to it only through the
Mumford--Tate conjecture, which is itself open. The two known inclusions, $\mathrm{algebraic} \subseteq \mathrm{Hodge}$ and $\mathrm{algebraic} \subseteq \mathrm{Tate}$, are joined by Deligne's theorem, which gives
$\mathrm{Hodge} \subseteq \mathrm{Tate}$ for abelian varieties (Corollary~\ref{absolute-hodge:corollary-tate-classes} and Theorem~\ref{absolute-hodge:theorem-deligne}), and Theorem~\ref{absolute-hodge:theorem-deligne-inclusion}
is the group-theoretic form of the same statement. For abelian varieties Faltings' theorem proves the Tate conjecture for divisors, and the Lefschetz $(1,1)$-theorem the Hodge conjecture for
divisors; in codimension $\geq 2$ both are open.
\end{remark}

\section{Exercises on The Tate Conjecture}\label{absolute-hodge:section-exercises-tate-conjecture}

\begin{exercise}\label{absolute-hodge:exercise-frobenius}
Let $k = \mathbb{F}_q$ and let $F$ be the geometric Frobenius, acting on $H^{2p}(\bar X, \QQ_\ell)$. Assuming that $F$ acts semisimply, show that $\dim T^p_\ell(\bar X)$ is the number of eigenvalues of $F$ on
$H^{2p}(\bar X, \QQ_\ell)$, counted with multiplicity, of the form $\zeta q^p$ with $\zeta$ a root of unity.
\end{exercise}

\begin{solution}
$G_{\mathbb{F}_{q^m}}$ is topologically generated by $F^m$, and the geometric Frobenius acts on $\QQ_\ell(1)$ by $q^{-1}$. So a class in $H^{2p}(\bar X, \QQ_\ell(p)) = H^{2p}(\bar X, \QQ_\ell) \otimes \QQ_\ell(p)$ is fixed by
$G_{\mathbb{F}_{q^m}}$ if and only if, as a class in $H^{2p}(\bar X, \QQ_\ell)$, it satisfies $F^mv = q^{pm}v$. Since $F$ is semisimple, the union over $m$ of the eigenspaces of $F^m$ for $q^{pm}$ is the sum of the
eigenspaces of $F$ for the eigenvalues $\lambda$ with $\lambda^m = q^{pm}$ for some $m$, that is, $\lambda = \zeta q^p$. The union is an increasing union of subspaces, so it is the largest of them.
\end{solution}

\begin{exercise}\label{absolute-hodge:exercise-elliptic-mt}
Let $E$ be an elliptic curve over a number field $k \subseteq \CC$ with complex multiplication by an imaginary quadratic field $K$, all of whose endomorphisms are defined over $k$. Show that $G_\ell^0$ is contained in the
torus $(\mathrm{Res}_{K/\QQ}\mathbb{G}_m)_{\QQ_\ell}$, and deduce that $G_\ell^0 \subseteq \mathrm{MT}(H^1(E_\CC))_{\QQ_\ell}$ directly, without Theorem~\ref{absolute-hodge:theorem-deligne-inclusion}.
\end{exercise}

\begin{solution}
The action of $K = \End(E) \otimes \QQ$ on $V_\ell E$ commutes with $G_k$, since the endomorphisms are defined over $k$. $V_\ell E$ is free of rank one over $K \otimes \QQ_\ell$, so the commutant of $K \otimes \QQ_\ell$
in $\End(V_\ell E)$ is $K \otimes \QQ_\ell$ itself. Hence $\rho_\ell(G_k) \subseteq (K \otimes \QQ_\ell)^\times = (\mathrm{Res}_{K/\QQ}\mathbb{G}_m)(\QQ_\ell)$, and the same holds for the Zariski closure. By
Proposition~\ref{hodge-structures:proposition-elliptic-curves}, this torus is $\mathrm{MT}(H^1(E_\CC))$ (the same argument applies to the dual $H^1$).
\end{solution}

\begin{exercise}\label{absolute-hodge:exercise-supersingular}
Let $E$ be a supersingular elliptic curve over $\mathbb{F}_p$, and $k = \mathbb{F}_{p^2}$ a field over which all endomorphisms of $E$ are defined, so that $\End(E_k) \otimes \QQ$ is a quaternion algebra of rank $4$. Using
Theorem~\ref{absolute-hodge:theorem-isogeny}, show that $\rho(E \times E) = 6$ over $k$, and that all of $H^2(\overline{E \times E}, \QQ_\ell(1))$ consists of Tate classes over $k$. Why is this impossible over $\CC$?
\end{exercise}

\begin{solution}
By Example~\ref{absolute-hodge:example-products-curves}, the N\'eron--Severi group of $E \times E$ has rank $2 + \operatorname{rank}\Hom(E, E) = 2 + 4 = 6 = b_2(E \times E)$. By Corollary~\ref{absolute-hodge:corollary-divisors-abelian} the
Tate classes are spanned by divisor classes, and conversely divisor classes over $k$ are Tate; since the divisor classes span a $6$-dimensional space, every class in $H^2$ is a Tate class. Over $\CC$ the rank of
$\End(E)$ is at most $2$ (Proposition~\ref{algebraic-groups:proposition-elliptic-endomorphisms}), so $\rho(E \times E) \leq 4 < 6$; equivalently $h^{2,0}(E \times E) = 1$ prevents all of $H^2$ from being of type $(1,1)$.
\end{solution}

\section{The Standard Conjectures and Motives}\label{absolute-hodge:section-standard-conjectures}
In the 1960s Grothendieck proposed a small set of conjectures about algebraic cycles, the \emph{standard conjectures}, designed to prove the Weil conjectures and to build a category of
\emph{motives}: a universal cohomology theory through which every Weil cohomology factors. Deligne proved the Weil conjectures by other means, and the standard conjectures remain open. They
are much weaker than the Hodge conjecture, which implies them in characteristic $0$ (Theorem~\ref{absolute-hodge:theorem-hodge-implies-standard}), yet even the Lefschetz standard
conjecture, which asks for one explicit class to be algebraic, is open in general. The following sections state the conjectures, proves the implications among them, constructs pure motives, proves Jannsen's
semisimplicity theorem, reformulates the Hodge conjecture motivically, and describes the conjectures of Bloch and Beilinson on the kernel of the cycle class map, the other half of the theory of
algebraic cycles. References: \cite{grothendieck-standard}, \cite{kleiman-1968}, \cite{kleiman-1994}, \cite{scholl-1994}, \cite{andre-motifs}, \cite{jannsen-1994}, and the source
\emph{The Hodge Conjecture} (Chapters 8 and 9).

\noindent\textbf{Prerequisites.} Chapter~\ref{intersection-theory} (Algebraic Cycles and Chow Groups), Section~\ref{kahler:section-lefschetz-theorems} (Hard Lefschetz and the Hodge--Riemann Relations) and
Section~\ref{absolute-hodge:section-tate-conjecture} (The Tate Conjecture).

\noindent\textbf{Conventions.} $k$ is a field, and varieties are nonsingular projective over $k$; $X$ has pure dimension $n$ unless stated otherwise. We fix a Weil cohomology theory $H^\bullet$ with
coefficients in a field $K$ of characteristic $0$ (Definition~\ref{etale-cohomology:definition-weil-cohomology}): $\ell$-adic cohomology of $X \times_k\bar k$ in general, or Betti cohomology of $X_\CC$ with
$K = \QQ$ when $k \subseteq \CC$. Tate twists are suppressed from the notation. For $\sim$ one of rational, homological (with respect to $H^\bullet$) or numerical equivalence
(Definition~\ref{intersection-theory:definition-equivalences}) we write
\[
A^p_\sim(X) = (Z^p(X) \otimes_\ZZ\QQ)/\!\sim, \qquad A^\bullet_\sim(X) = \bigoplus_pA^p_\sim(X),
\]
so that $A^p_{\mathrm{rat}}(X) = \CH^p(X)_\QQ$ and $A^p_{\mathrm{hom}}(X)$ is the $\QQ$-span of the image of the cycle class map in $H^{2p}(X)$, whose elements we call \emph{algebraic classes}. We fix an ample divisor on
$X$ with class $h \in A^1_{\mathrm{hom}}(X)$ and write $L = h \cup - \colon H^i(X) \to H^{i+2}(X)$.

\section{Correspondences in cohomology}\label{absolute-hodge:section-correspondences}

\begin{proposition}\label{absolute-hodge:proposition-cohomological-correspondences}
Let $X$, $Y$, $W$ be varieties, $X$ of dimension $n$.
\begin{enumerate}
\item A class $\gamma \in H^{2n+2r}(X \times Y)$ acts by $\gamma_*(x) = \mathrm{pr}_{2*}(\mathrm{pr}_1^*x \cup \gamma)$, which maps $H^i(X)$ to $H^{i+2r}(Y)$, and the K\"unneth formula together with Poincar\'e duality
give an isomorphism $H^{2n+2r}(X \times Y) \cong \bigoplus_i\Hom_K(H^i(X), H^{i+2r}(Y))$, $\gamma \mapsto (\gamma_*|_{H^i(X)})_i$.
\item For $\gamma \in H^\bullet(X \times Y)$ and $\delta \in H^\bullet(Y \times W)$, the class $\delta \circ \gamma = \mathrm{pr}_{13*}(\mathrm{pr}_{12}^*\gamma \cup \mathrm{pr}_{23}^*\delta)$ satisfies
$(\delta \circ \gamma)_* = \delta_* \circ \gamma_*$, the class of the diagonal acts as the identity, and $\mathrm{cl}(\Gamma' \circ \Gamma) = \mathrm{cl}(\Gamma') \circ \mathrm{cl}(\Gamma)$ for algebraic correspondences
(Definition~\ref{intersection-theory:definition-correspondence}).
\item (Lefschetz trace formula) For $\alpha, \beta \in H^{2n}(X \times X)$, $\int_{X \times X}\alpha \cup {}^t\beta = \sum_i(-1)^i\operatorname{tr}\big((\beta \circ \alpha)_* \mid H^i(X)\big)$, where ${}^t\beta$ is the pullback of
$\beta$ under the exchange of the factors. The right side is the \emph{supertrace} $\operatorname{str}((\beta \circ \alpha)_*)$.
\end{enumerate}
\end{proposition}

These are formal consequences of the axioms of a Weil cohomology \cite{kleiman-1968}; (2) is the cohomological version of Proposition~\ref{intersection-theory:proposition-correspondences}, and (3)
for $\beta = \Delta_X$ is the Lefschetz trace formula of Theorem~\ref{etale-cohomology:theorem-l-adic-properties}(5). By (2), an algebraic correspondence maps algebraic classes to algebraic classes:
$\mathrm{cl}(\Gamma)_*\mathrm{cl}(Z) = \mathrm{cl}(\Gamma_*Z)$. We call a $K$-linear endomorphism of $H^\bullet(X)$ \emph{algebraic} if it is $\gamma_*$ for $\gamma$ in the $\QQ$-span of classes of algebraic cycles on
$X \times X$. Sums and composites of algebraic endomorphisms are algebraic, and $L = (\Delta_*h)_*$ is algebraic.

\begin{lemma}\label{absolute-hodge:lemma-finite-dimensional}
$A^p_{\mathrm{num}}(X)$ is a finite-dimensional $\QQ$-vector space of dimension at most $\dim_KH^{2p}(X)$, and the pairing $A^p_{\mathrm{num}}(X) \times A^{n-p}_{\mathrm{num}}(X) \to \QQ$,
$(\alpha, \beta) \mapsto \deg(\alpha \cdot \beta)$, is perfect. If $K = \QQ$, then $A^p_{\mathrm{hom}}(X)$ is finite-dimensional too.
\end{lemma}

\begin{proof}
Homological equivalence implies numerical equivalence (Proposition~\ref{intersection-theory:proposition-equivalences}), and $\deg(\alpha \cdot \beta) = \int\mathrm{cl}(\alpha) \cup \mathrm{cl}(\beta)$. Let
$\alpha_1, \ldots, \alpha_N \in A^p_{\mathrm{num}}(X)$ be linearly independent. The map $A^{n-p}_{\mathrm{num}}(X) \to \QQ^N$, $\beta \mapsto (\deg(\alpha_i \cdot \beta))_i$, is surjective: otherwise a nonzero linear form on $\QQ^N$
vanishes on its image, which gives a nontrivial combination of the $\alpha_i$ that pairs to zero with every cycle, that is, is zero in $A^p_{\mathrm{num}}$. Choose $\beta_j$ with $\deg(\alpha_i \cdot \beta_j) = \delta_{ij}$. If
$\sum_ic_i\mathrm{cl}(\alpha_i) = 0$ with $c_i \in K$, cupping with $\mathrm{cl}(\beta_j)$ and integrating gives $c_j = 0$. So the $\mathrm{cl}(\alpha_i)$ are $K$-linearly independent in $H^{2p}(X)$ and
$N \leq \dim_KH^{2p}(X)$. The pairing is nondegenerate by definition, hence perfect between finite-dimensional spaces. If $K = \QQ$, $A^p_{\mathrm{hom}}(X)$ is a subspace of $H^{2p}(X)$.
\end{proof}

\begin{remark}\label{absolute-hodge:remark-hom-finite}
For $\ell$-adic cohomology, $A^p_{\mathrm{hom}}(X)$ is a $\QQ$-subspace of a $\QQ_\ell$-vector space, and it is not known in general that it is finite-dimensional over $\QQ$, or that
$A^p_{\mathrm{hom}}(X) \otimes_\QQ\QQ_\ell \to H^{2p}(X)$ is injective. Both follow from the equality of homological and numerical equivalence (Conjecture~\ref{absolute-hodge:conjecture-d} below), by the
proof of Lemma~\ref{absolute-hodge:lemma-finite-dimensional}.
\end{remark}

\section{The standard conjectures}\label{absolute-hodge:section-conjectures}

By hard Lefschetz (an axiom of Weil cohomologies; for $\ell$-adic cohomology a theorem of Deligne, Theorem~\ref{etale-cohomology:theorem-l-adic-properties}(6)), $L^{n-i} \colon H^i(X) \to H^{2n-i}(X)$ is an
isomorphism for $i \leq n$. As in Section~\ref{kahler:section-lefschetz-theorems}, this gives the Lefschetz decomposition $H^i = \bigoplus_jL^jP^{i-2j}$, with $P^j = \ker(L^{n-j+1} \colon H^j \to H^{2n-j+2})$ for $j \leq n$, and a unique
operator $\Lambda$ with $[L, \Lambda] = h_{\deg}$ and $[h_{\deg}, \Lambda] = -2\Lambda$, where $h_{\deg}$ acts on $H^i$ by $i - n$ (Proposition~\ref{kahler:proposition-lambda-unique}); explicitly
$\Lambda(L^rx) = r(n - j - r + 1)L^{r-1}x$ for $x \in P^j$. So $H^\bullet(X)$ is a representation of $\mathfrak{sl}_2$ with $e = L$, $f = \Lambda$, $h = h_{\deg}$, in which $H^i$ is the weight space of weight $i - n$ and
$P^j$ consists of lowest weight vectors of weight $j - n$.

\begin{conjecture}[Lefschetz standard conjecture $B(X)$]\label{absolute-hodge:conjecture-b}
The operator $\Lambda$ is algebraic.
\end{conjecture}

\begin{conjecture}[K\"unneth standard conjecture $C(X)$]\label{absolute-hodge:conjecture-c}
The K\"unneth components $\pi_i \in H^{2n-i}(X) \otimes H^i(X) \subseteq H^{2n}(X \times X)$ of the class of the diagonal are algebraic; equivalently, the projectors of $H^\bullet(X)$ onto the summands $H^i(X)$
are algebraic.
\end{conjecture}

\begin{conjecture}[Standard conjecture $D(X)$]\label{absolute-hodge:conjecture-d}
Homological and numerical equivalence coincide on $Z^\bullet(X) \otimes \QQ$: $A^\bullet_{\mathrm{hom}}(X) = A^\bullet_{\mathrm{num}}(X)$.
\end{conjecture}

\begin{conjecture}[Hodge standard conjecture $\mathrm{Hdg}(X)$]\label{absolute-hodge:conjecture-hdg}
For $p \leq n/2$, the symmetric bilinear form $(x, y) \mapsto (-1)^p\int_XL^{n-2p}(x \cup y)$ is positive definite on the $\QQ$-space $A^p_{\mathrm{hom}}(X) \cap P^{2p}(X)$ of primitive algebraic classes.
\end{conjecture}

The equivalence in $C(X)$ is Proposition~\ref{absolute-hodge:proposition-cohomological-correspondences}(1): $(\pi_i)_*$ is the projector onto $H^i(X)$.

\begin{lemma}\label{absolute-hodge:lemma-sl2-polynomials}
Let $V$ be a finite-dimensional representation of $\mathfrak{sl}_2$ over a field of characteristic $0$, with weight spaces $V_w$ for $h$.
\begin{enumerate}
\item The projector of $V$ onto $V_w$, along the other weight spaces, is a polynomial in $h$.
\item The projector onto the isotypic component of $V(m)$ is a polynomial in the Casimir element $c = ef + fe + \frac12h^2$, which acts on $V(m)$ by $\frac12m(m + 2)$.
\item For $j \geq 0$, $e^j \colon V_{-j} \to V_j$ is bijective, and its inverse is the restriction to $V_j$ of $P(f^je^j)f^j$ for a polynomial $P$.
\item If $v$ is a lowest weight vector of weight $-m$ and $0 \leq j \leq m$, then $f^je^jv = c_{m,j}v$ with $c_{m,j}$ a nonzero integer.
\end{enumerate}
\end{lemma}

\begin{proof}
We use the basis $v_0, \ldots, v_m$ of $V(m)$ from Theorem~\ref{representation-theory:theorem-sl2-irreducibles}, with $hv_i = (m - 2i)v_i$, $fv_i = (i + 1)v_{i+1}$, $ev_i = (m - i + 1)v_{i-1}$, and the decomposition of $V$
into irreducibles. (1) $h$ is diagonalisable with finitely many eigenvalues $w_1, \ldots, w_s$, and the projector onto $V_{w_a}$ is $\prod_{b \neq a}(h - w_b)/(w_a - w_b)$. (2) On $v_0$, $ev_0 = 0$ gives
$efv_0 = [e, f]v_0 = hv_0 = mv_0$ and $fev_0 = 0$, so $cv_0 = (m + \frac12m^2)v_0$; as $c$ commutes with $e$, $f$, $h$, it acts on $V(m)$ by this scalar. The scalars $\frac12m(m + 2)$ are distinct for distinct
$m \geq 0$, so the projectors are Lagrange polynomials in $c$. (4) The lowest weight vector of $V(m)$ is $v_m$. Applying $e^j$ multiplies it by $\prod_{s=0}^{j-1}(s + 1) = j!$ and gives $v_{m-j}$; applying $f^j$ multiplies by
$\prod_{t=0}^{j-1}(m - j + t + 1)$ and returns to $v_m$. All factors are positive. (3) The weight $-j$ vector of $V(m)$, when it exists ($m \geq j$, $m \equiv j \bmod 2$), is $v_i$ with $i = (m + j)/2$; applying $e^j$ and
then $f^j$ multiplies it by $\prod_{s=0}^{j-1}(m - i + s + 1)\prod_{t=0}^{j-1}(i - j + t + 1)$, a product of integers $\geq (m - j)/2 + 1 \geq 1$. So $f^je^j$ acts on $V_{-j}$ diagonalisably with finitely many nonzero eigenvalues
$c_1, \ldots, c_r$, and $P(f^je^j)$ inverts it for the Lagrange polynomial with $P(c_a) = c_a^{-1}$. Then $P(f^je^j)f^j \circ e^j$ is the identity of $V_{-j}$, so $e^j$ is injective there; as $\dim V_{-j} = \dim V_j$, it is
bijective, with inverse $P(f^je^j)f^j|_{V_j}$.
\end{proof}

\begin{proposition}\label{absolute-hodge:proposition-b-implies}
Assume $B(X)$. Then:
\begin{enumerate}
\item $C(X)$ holds;
\item for $p \leq n/2$, $L^{n-2p} \colon A^p_{\mathrm{hom}}(X) \to A^{n-p}_{\mathrm{hom}}(X)$ is bijective (\emph{Grothendieck's conjecture $A(X, L)$});
\item if $x = \sum_jL^jx_j$ is the Lefschetz decomposition of an algebraic class, with $x_j$ primitive, then every $x_j$ is algebraic.
\end{enumerate}
\end{proposition}

\begin{proof}
By $B(X)$, $L$ and $\Lambda$ are algebraic, hence so is every noncommutative polynomial in them, in particular $h_{\deg} = [L, \Lambda]$ and the Casimir element. (1) By
Lemma~\ref{absolute-hodge:lemma-sl2-polynomials}(1), the projector onto $H^i$ is a polynomial in $h_{\deg}$. (2) $L^{n-2p}$ maps algebraic classes to algebraic classes and is injective on $H^{2p}$ by hard
Lefschetz. By Lemma~\ref{absolute-hodge:lemma-sl2-polynomials}(3), its inverse $H^{2n-2p} \to H^{2p}$ is the restriction of an algebraic endomorphism, so it maps algebraic classes to algebraic classes.
(3) The summand $L^jP^{i-2j}$ of $H^i$ is the intersection of $H^i$ with the isotypic component of $V(n - i + 2j)$, since a primitive class in $P^{i-2j}$ is a lowest weight vector of weight $i - 2j - n$. By
Lemma~\ref{absolute-hodge:lemma-sl2-polynomials}(1),(2) the projector onto $L^jP^{i-2j}$ is algebraic, so $L^jx_j$ is algebraic. By Lemma~\ref{absolute-hodge:lemma-sl2-polynomials}(4),
$x_j = c^{-1}\Lambda^jL^jx_j$ for a nonzero integer $c$, so $x_j$ is algebraic.
\end{proof}

\begin{proposition}\label{absolute-hodge:proposition-hdg-char-zero}
If $k \subseteq \CC$ and $H^\bullet$ is Betti cohomology, then $\mathrm{Hdg}(X)$ holds.
\end{proposition}

\begin{proof}
Algebraic classes in $H^{2p}(X_\CC; \QQ)$ are rational classes of type $(p,p)$ (Corollary~\ref{intersection-theory:corollary-hodge-classes}). The class $h$ is represented by a K\"ahler form $\omega$, the restriction of a
multiple of the Fubini--Study form. For a nonzero primitive rational class $x$ of type $(p,p)$, the Hodge--Riemann relation (Theorem~\ref{kahler:theorem-hodge-riemann}(4)) with $k = 2p$ and $\bar x = x$ reads
$Q(x, x) > 0$, where $Q(x, x) = (-1)^{2p(2p-1)/2}\int x \cup x \cup [\omega]^{n-2p} = (-1)^p\int L^{n-2p}(x \cup x)$, since $p(2p - 1) \equiv p \pmod 2$.
\end{proof}

\begin{proposition}\label{absolute-hodge:proposition-b-hdg-d}
$B(X)$ and $\mathrm{Hdg}(X)$ together imply $D(X)$.
\end{proposition}

\begin{proof}
First let $p \leq n/2$, and put $\langle x, y\rangle = \int L^{n-2p}(x \cup y)$ on $H^{2p}$. The Lefschetz decomposition is orthogonal for $\langle\ ,\ \rangle$: for $a \in P^{2p-2j}$, $b \in P^{2p-2j'}$ and $j < j'$,
$\langle L^ja, L^{j'}b\rangle = \int b \cup L^{n-2p+j+j'}a$ vanishes because $n - 2p + j + j' \geq n - (2p - 2j) + 1$, so $L^{n-2p+j+j'}a = 0$ by primitivity. Let $x \in A^p_{\mathrm{hom}}(X)$ with $\langle x, y\rangle = 0$ for every
$y \in A^p_{\mathrm{hom}}(X)$, and write $x = \sum_jL^jx_j$. By Proposition~\ref{absolute-hodge:proposition-b-implies}(3) the $x_j$ are algebraic, so $y = L^jx_j$ is allowed, and
$0 = \langle x, L^jx_j\rangle = \langle L^jx_j, L^jx_j\rangle = \int L^{n-2(p-j)}(x_j \cup x_j)$. By $\mathrm{Hdg}(X)$ in codimension $p - j$, $x_j = 0$. Hence $\langle\ ,\ \rangle$ is nondegenerate on $A^p_{\mathrm{hom}}(X)$.
Now if $x \in A^p_{\mathrm{hom}}(X)$ is numerically trivial, then $\langle x, y\rangle = \deg(x \cdot h^{n-2p} \cdot y) = 0$ for all algebraic $y$, so $x = 0$.

Now let $p > n/2$ and let $x \in A^p_{\mathrm{hom}}(X)$ be numerically trivial. By Proposition~\ref{absolute-hodge:proposition-b-implies}(2) applied to $n - p < n/2$, $L^{2p-n} \colon A^{n-p}_{\mathrm{hom}} \to A^p_{\mathrm{hom}}$ is
bijective, so $x = L^{2p-n}y$ with $y$ algebraic, and every algebraic class of codimension $p$ is $L^{2p-n}u$ with $u \in A^{n-p}_{\mathrm{hom}}$. Then $\deg(y \cdot L^{2p-n}u) = \deg(x \cdot u) = 0$ for all such $u$, so $y$ is
numerically trivial, hence $y = 0$ by the first case, and $x = 0$.
\end{proof}

\begin{theorem}\label{absolute-hodge:theorem-hodge-implies-standard}
Let $k \subseteq \CC$ and let $H^\bullet$ be Betti cohomology. If the Hodge conjecture holds for $X \times X$, then $B(X)$, $C(X)$ and $D(X)$ hold.
\end{theorem}

\begin{proof}
The operator $\Lambda$ is defined over $\QQ$ and is a morphism of Hodge structures of type $(-1,-1)$ (see the proof of Corollary~\ref{absolute-hodge:corollary-lambda}), so its class in
$H^{2n-2}(X \times X; \QQ)$ is a Hodge class (Proposition~\ref{hodge-structures:proposition-hodge-classes-morphisms}). By the Hodge conjecture it is algebraic, which is $B(X)$. Then $C(X)$ follows from
Proposition~\ref{absolute-hodge:proposition-b-implies}, and $D(X)$ from Propositions~\ref{absolute-hodge:proposition-hdg-char-zero} and~\ref{absolute-hodge:proposition-b-hdg-d}.
\end{proof}

\begin{remark}\label{absolute-hodge:remark-grothendieck-plan}
Grothendieck showed that the standard conjectures for $X$ and $X \times X$ over a finite field imply the Riemann hypothesis part of the Weil conjectures, by an argument modelled on Serre's K\"ahler analogue:
the graph of Frobenius is an algebraic correspondence, and the positivity of $\mathrm{Hdg}(X \times X)$ bounds its eigenvalues \cite{grothendieck-standard}, \cite{kleiman-1968}. Deligne's proof
\cite{deligne-weil-1} avoided the standard conjectures. Katz and Messing then deduced $C(X)$ over finite fields from Deligne's theorem \cite{katz-messing-1974}. In positive characteristic,
$\mathrm{Hdg}(X)$ is known in general only for surfaces, where it is the Hodge index theorem \cite[Theorem V.1.9]{hartshorne-ag}.
\end{remark}

\begin{theorem}[Known cases]\label{absolute-hodge:theorem-known-standard-conjectures}
\begin{enumerate}
\item $B(X)$ holds if $\dim X \leq 2$, if $X$ is an abelian variety \cite{lieberman-1968}, for products of varieties for which it holds \cite{kleiman-1968}, for varieties whose cohomology is spanned by algebraic
classes (projective spaces, Grassmannians, flag varieties), for smooth hypersurfaces and complete intersections in projective space (Exercise~\ref{absolute-hodge:exercise-hypersurface}), and for hyper-K\"ahler varieties
deformation equivalent to Hilbert schemes of points on K3 surfaces \cite{charles-markman-2013}.
\item $C(X)$ holds over finite fields \cite{katz-messing-1974}, and whenever $B(X)$ holds.
\item $D(X)$ holds for divisors in any characteristic, and in characteristic $0$ for abelian varieties \cite{lieberman-1968}.
\item $\mathrm{Hdg}(X)$ holds in characteristic $0$ (Proposition~\ref{absolute-hodge:proposition-hdg-char-zero}) and for surfaces in any characteristic.
\end{enumerate}
\end{theorem}

See \cite{kleiman-1994} for a survey and for the classical cases. For a surface, the components of $\Lambda$ on $H^2 \to H^0$ and $H^4 \to H^2$ are expressed through $h$ and the class of a point, and
the component $H^3 \to H^1$, the inverse of $L \colon H^1 \to H^3$, is induced by a divisorial correspondence built from the Picard and Albanese varieties. Milne showed that $\mathrm{Hdg}$ for abelian
varieties in positive characteristic would follow from the Hodge conjecture for complex abelian varieties of CM type \cite{milne-2002}.

\begin{remark}\label{absolute-hodge:remark-weakness}
$B(X)$ asks for the algebraicity of one specific Hodge class on $X \times X$, while the Hodge conjecture asks for all of them, and still $B(X)$ is open for most varieties of dimension $\geq 3$. For a
threefold, $B(X)$ essentially amounts to the algebraicity of the inverse of $L \colon H^2(X) \to H^4(X)$, a class in $H^2(X) \otimes H^2(X)$ whose transcendental part is not accessible by divisors. The class of $\Lambda$ is absolute
Hodge (Corollary~\ref{absolute-hodge:corollary-lambda}), but no algebraic cycle representing it is known in general.
\end{remark}

\section{Pure motives}\label{absolute-hodge:section-motives}

\begin{definition}\label{absolute-hodge:definition-motives}
Let $\sim$ be rational, homological or numerical equivalence. The category $\mathcal{M}_\sim(k)$ of \emph{pure motives} over $k$ (with $\QQ$-coefficients) has as objects triples $(X, e, m)$ with $X$ a
nonsingular projective variety of pure dimension, $e \in A^{\dim X}_\sim(X \times X)$ an idempotent, $e \circ e = e$, and $m \in \ZZ$, together with their finite formal direct sums. Morphisms between triples are
\[
\Hom\big((X, e, m), (Y, f, m')\big) = f \circ A^{\dim X + m' - m}_\sim(X \times Y) \circ e,
\]
with composition of correspondences, and are extended to direct sums by matrices. The \emph{motive} of $X$ is $h(X) = (X, \Delta_X, 0)$, and a morphism $\phi \colon X \to Y$ gives
$h(\phi) = {}^t\Gamma_\phi \colon h(Y) \to h(X)$. The \emph{unit motive} is $\mathbf{1} = h(\Spec k)$, the \emph{Lefschetz motive} is $\mathbb{L} = (\Spec k, \mathrm{id}, -1)$, and the \emph{Tate twist} is
$(X, e, m)(r) = (X, e, m + r)$.
\end{definition}

Composition is well defined because $\sim$ is compatible with pushforward, pullback and intersection product (Proposition~\ref{intersection-theory:proposition-correspondences}). The category
$\mathcal{M}_\sim(k)$ is $\QQ$-linear and additive, and \emph{pseudo-abelian}: an idempotent endomorphism $\phi = f \circ \phi \circ f$ of $(X, f, m)$ has image $(X, \phi, m)$. The product
$(X, e, m) \otimes (Y, f, m') = (X \times Y, e \times f, m + m')$ makes it a tensor category, with $(X, e, m) \otimes \mathbb{L} = (X, e, m - 1)$ \cite{scholl-1994}.

\begin{example}\label{absolute-hodge:example-projective-line}
In $\mathcal{M}_{\mathrm{rat}}(k)$, $h(\PP^1) \cong \mathbf{1} \oplus \mathbb{L}$. Let $x$ be a $k$-point of $\PP^1$, and put $e_0 = [x \times \PP^1]$ and $e_2 = [\PP^1 \times x]$ in $\CH^1(\PP^1 \times \PP^1)$. The diagonal has
bidegree $(1,1)$ and $\CH^1(\PP^1 \times \PP^1)$ is free on $e_0$ and $e_2$, so $\Delta = e_0 + e_2$. By the composition formula of Definition~\ref{intersection-theory:definition-correspondence}, on
$\PP^1 \times \PP^1 \times \PP^1$:
\begin{enumerate}
\item $e_0 \circ e_0 = \mathrm{pr}_{13*}([x \times \PP^1 \times \PP^1] \cdot [\PP^1 \times x \times \PP^1]) = \mathrm{pr}_{13*}[x \times x \times \PP^1] = e_0$;
\item $e_2 \circ e_2 = \mathrm{pr}_{13*}([\PP^1 \times x \times \PP^1] \cdot [\PP^1 \times \PP^1 \times x]) = \mathrm{pr}_{13*}[\PP^1 \times x \times x] = e_2$;
\item $e_2 \circ e_0 = \mathrm{pr}_{13*}([x \times \PP^1 \times \PP^1] \cdot [\PP^1 \times \PP^1 \times x]) = \mathrm{pr}_{13*}[x \times \PP^1 \times x] = 0$, since $\mathrm{pr}_{13}$ contracts the middle factor;
\item $e_0 \circ e_2 = \mathrm{pr}_{13*}([\PP^1 \times x \times \PP^1]^2) = 0$, since $[x]^2 = 0$ in $\CH^\bullet(\PP^1)$.
\end{enumerate}
So $e_0$, $e_2$ are orthogonal idempotents with sum $\Delta$, and $h(\PP^1) = (\PP^1, e_0, 0) \oplus (\PP^1, e_2, 0)$. The correspondences $a = [\PP^1] \in \CH^0(\Spec k \times \PP^1)$ and $b = [x] \in \CH^1(\PP^1 \times \Spec k)$
satisfy $b \circ a = \mathrm{id}_{\mathbf{1}}$ and $a \circ b = e_0$, so $(\PP^1, e_0, 0) \cong \mathbf{1}$. Similarly $a' = [x] \in \CH^1(\Spec k \times \PP^1)$, a morphism $\mathbb{L} \to h(\PP^1)$, and $b' = [\PP^1] \in \CH^0(\PP^1 \times \Spec k)$,
a morphism $h(\PP^1) \to \mathbb{L}$, satisfy $b' \circ a' = \mathrm{id}_{\mathbb{L}}$ and $a' \circ b' = e_2$, so $(\PP^1, e_2, 0) \cong \mathbb{L}$. In cohomology, $e_0$ and $e_2$ are the K\"unneth components $\pi_0$ and $\pi_2$.
\end{example}

\begin{theorem}[Jannsen]\label{absolute-hodge:theorem-jannsen}
$\mathcal{M}_{\mathrm{num}}(k)$ is a semisimple abelian category. Conversely, if $\sim$ is an adequate equivalence relation such that $\mathcal{M}_\sim(k)$ is semisimple abelian, then $\sim$ is numerical
equivalence.
\end{theorem}

\begin{proof}[Proof of the first statement]
\emph{Step 1: for every $X$ of dimension $n$, the algebra $E = A^n_{\mathrm{num}}(X \times X)$, with composition of correspondences, is semisimple.} It is finite-dimensional over $\QQ$ by
Lemma~\ref{absolute-hodge:lemma-finite-dimensional}, so its Jacobson radical $R$ is a nilpotent two-sided ideal. Let $\alpha \in Z^n(X \times X) \otimes \QQ$ have image in $R$, and let $\beta$ be any cycle of
dimension $n$ on $X \times X$. Put $\gamma = \beta \circ \alpha$. Its image lies in $R$, so the image of $\gamma^{\circ s}$ in $E$ is $0$ for some $s \geq 1$, and then so is the image of $\gamma^{\circ t}$ for all $t \geq s$.
By Proposition~\ref{absolute-hodge:proposition-cohomological-correspondences}(3) with ${}^t\Delta = \Delta$,
\[
\operatorname{str}((\gamma_*)^t) = \int\mathrm{cl}(\gamma^{\circ t}) \cup \mathrm{cl}(\Delta) = \deg(\gamma^{\circ t} \cdot \Delta) = 0 \qquad (t \geq s),
\]
because $\gamma^{\circ t}$ is numerically trivial. Let $\lambda_1, \ldots, \lambda_a$ and $\mu_1, \ldots, \mu_b$ be the eigenvalues, with multiplicity, of $\gamma_*$ on $H^{\mathrm{even}}$ and on $H^{\mathrm{odd}}$, in an algebraic closure of $K$.
Then $\sum_i\lambda_i^t - \sum_j\mu_j^t = 0$ for $t \geq s$, so the rational function
\[
R(z) = \sum_i\frac{\lambda_iz}{1 - \lambda_iz} - \sum_j\frac{\mu_jz}{1 - \mu_jz} = \sum_{t \geq 1}\Big(\sum_i\lambda_i^t - \sum_j\mu_j^t\Big)z^t
\]
is a polynomial. Its poles at the points $1/\lambda$, for nonzero eigenvalues $\lambda$, must therefore cancel, which means that every nonzero value occurs equally often among the $\lambda_i$ and the $\mu_j$. Hence
$\sum_i\lambda_i^t = \sum_j\mu_j^t$ for all $t \geq 1$, and in particular $\operatorname{str}(\gamma_*) = 0$. By the trace formula again, $\deg(\alpha \cdot {}^t\beta) = \operatorname{str}((\beta \circ \alpha)_*) = 0$. As
$\beta$ was arbitrary and every cycle of dimension $n$ on $X \times X$ is of the form ${}^t\beta$, $\alpha$ is numerically trivial. So $R = 0$, and $E$ is semisimple.

\emph{Step 2: every object is a direct summand of some $h(Y)(m)$.} The triple $(X, e, m)$ is the image of the idempotent $e$ of $h(X)(m)$. Given finitely many triples $(X_i, e_i, m_i)$, choose $m \geq \max m_i$. Then
$(X_i, e_i, m_i) = (X_i, e_i, m) \otimes \mathbb{L}^{\otimes(m - m_i)}$, and $\mathbb{L}^{\otimes r}$ is a summand of $h(\PP^1)^{\otimes r} = h((\PP^1)^r)$ by Example~\ref{absolute-hodge:example-projective-line} (which
holds modulo numerical equivalence as well). So each $(X_i, e_i, m_i)$ is a summand of $h(X_i \times (\PP^1)^{m-m_i})(m)$, and their direct sum is a summand of $h(Y)(m)$ with $Y$ the disjoint union of the
$X_i \times (\PP^1)^{m-m_i}$.

\emph{Step 3: conclusion.} By Step 2, the endomorphism algebra of an object $M$ is $\varepsilon E\varepsilon$ for an idempotent $\varepsilon$ of $E = \End(h(Y)(m)) = A^{\dim Y}_{\mathrm{num}}(Y \times Y)$ (for disconnected $Y$,
a product of such algebras over pairs of components, with the same argument). A corner $\varepsilon E\varepsilon$ of a finite-dimensional semisimple algebra is semisimple. Finally, a $\QQ$-linear additive pseudo-abelian
category in which every object has a finite-dimensional semisimple endomorphism algebra is abelian and semisimple \cite{jannsen-1992}.
\end{proof}

The converse statement is also in \cite{jannsen-1992}. Jannsen's theorem was a surprise: Grothendieck expected semisimplicity to require the standard conjectures.

\begin{remark}\label{absolute-hodge:remark-tannakian}
What the standard conjectures are still needed for is the \emph{Tannakian} structure. The natural commutativity constraint on $\mathcal{M}_{\mathrm{num}}(k)$ gives $h(X)$ the wrong dimension (the
supertrace $\sum(-1)^ib_i$, which can be negative); to correct it one changes the sign of the constraint on the odd part, which requires the K\"unneth projectors to be algebraic modulo numerical
equivalence. With $C$ and $D$, $\mathcal{M}_{\mathrm{num}}(k) = \mathcal{M}_{\mathrm{hom}}(k)$ is a semisimple Tannakian category, the realisation $H^\bullet$ is a fibre functor, and $\mathcal{M}_{\mathrm{num}}(k)$ is the category of
representations of a pro-reductive group, the \emph{motivic Galois group} \cite{jannsen-1992}, \cite{andre-motifs}, \cite{deligne-milne}. Unconditional substitutes exist: Deligne's motives for absolute Hodge
cycles \cite{deligne-milne} and Andr\'e's motives for motivated cycles \cite{andre-1996-motifs} are Tannakian and contain the motives of abelian varieties.
\end{remark}

\begin{proposition}[The Hodge conjecture, motivically]\label{absolute-hodge:proposition-motivic-hodge}
Let $k = \CC$ and let $\sim$ be homological equivalence for Betti cohomology. The Hodge realisation, which sends $(X, e, m)$ to the graded polarisable $\QQ$-Hodge structure $e_*H^\bullet(X; \QQ)(m)$, with
$H^i(X)(m)$ placed in degree $i - 2m$, is a faithful functor from $\mathcal{M}_{\mathrm{hom}}(\CC)$ to graded polarisable $\QQ$-Hodge structures. It is full if and only if the Hodge conjecture holds for all
nonsingular complex projective varieties.
\end{proposition}

\begin{proof}
A morphism $h(X)(m) \to h(Y)(m')$ is an element $\gamma \in A^{n+r}_{\mathrm{hom}}(X \times Y)$ with $n = \dim X$, $r = m' - m$, and it acts by $\gamma_* \colon H^i(X)(m) \to H^{i+2r}(Y)(m')$, which preserves the grading.
Faithfulness is Proposition~\ref{absolute-hodge:proposition-cohomological-correspondences}(1): an element of $A_{\mathrm{hom}}$ is determined by its class, which is determined by its action. Under the
isomorphism of Proposition~\ref{absolute-hodge:proposition-cohomological-correspondences}(1), which is compatible with Hodge structures up to the Tate twist, the grading-preserving morphisms of Hodge
structures $\bigoplus_iH^i(X)(m) \to \bigoplus_iH^{i+2r}(Y)(m')$ correspond to the Hodge classes in $H^{2n+2r}(X \times Y; \QQ)$ (Proposition~\ref{hodge-structures:proposition-hodge-classes-morphisms}), and the
morphisms of motives to the algebraic classes. So the realisation is full on the objects $h(X)(m)$ if and only if $\mathrm{Hdg}^{n+r}(X \times Y) = Z^{n+r}_{\mathrm{alg}}(X \times Y)$ for all $X$, $Y$, $r$. Taking $X = \Spec\CC$
gives the Hodge conjecture for $Y$; conversely the Hodge conjecture for all varieties includes $X \times Y$. Fullness for general objects follows, since they are summands of the $h(X)(m)$ and morphisms between
summands are cut out by the idempotents.
\end{proof}

\begin{remark}\label{absolute-hodge:remark-motivic-galois}
Proposition~\ref{absolute-hodge:proposition-motivic-hodge} is the cleanest statement of what the Hodge conjecture is \emph{for}: the Hodge structure of a variety would determine its motive, and every
Hodge-theoretic relation between varieties would be realised by correspondences. Under the Hodge conjecture, the standard conjectures hold (Theorem~\ref{absolute-hodge:theorem-hodge-implies-standard}),
$\mathcal{M}_{\mathrm{hom}}(\CC)$ is Tannakian, and the motivic Galois group of the Tannakian subcategory generated by $h(X)$ is the Mumford--Tate group of $H^\bullet(X; \QQ)$
(Section~\ref{hodge-structures:section-mumford-tate}), an identification of a geometric object with a Hodge-theoretic one. If the Hodge conjecture were false, the motives for absolute Hodge cycles
(Remark~\ref{absolute-hodge:remark-tannakian}) would survive, but they would no longer be built out of subvarieties.
\end{remark}

\section{The Bloch--Beilinson conjectures}\label{absolute-hodge:section-bloch-beilinson}

The Hodge conjecture concerns the image of the cycle class map. Its kernel, $\CH^p(X)_{\mathrm{hom}} = \ker(\mathrm{cl})$, is governed by a separate set of conjectures.

\begin{theorem}[Mumford]\label{absolute-hodge:theorem-mumford}
Let $S$ be a nonsingular complex projective surface with $h^{2,0}(S) > 0$. Then the group $\CH_0(S)_{\deg 0}$ of $0$-cycles of degree $0$ modulo rational equivalence is \emph{infinite-dimensional}: for no
variety $T$ and cycle $Z$ on $T \times S$ is the map $T \to \CH_0(S)_{\deg 0}$, $t \mapsto Z_t - Z_{t_0}$, surjective.
\end{theorem}

This is \cite{mumford-1968}; compare Example~\ref{intersection-theory:example-chow-zero-surfaces}. The proof shows that a surjective family would force every holomorphic $2$-form to vanish, by pulling
it back along the family. Chow groups thus see Hodge-theoretic information that cohomology classes of cycles do not.

\begin{conjecture}[Bloch]\label{absolute-hodge:conjecture-bloch}
If $S$ is a nonsingular complex projective surface with $h^{2,0}(S) = 0$, then the Albanese map $\CH_0(S)_{\deg 0} \to \mathrm{Alb}(S)$ is an isomorphism.
\end{conjecture}

Bloch's conjecture is known for surfaces not of general type \cite{bloch-kas-lieberman-1976} and in scattered cases of general type, and open in general \cite{bloch-lectures}. By the Bloch--Srinivas method
(Theorem~\ref{known-cases:theorem-bloch-srinivas}), its conclusion implies a decomposition of the diagonal, and conversely the conclusion of Mumford's theorem is the Hodge-theoretic obstruction to one.

\begin{conjecture}[Bloch--Beilinson filtration]\label{absolute-hodge:conjecture-bloch-beilinson}
For every nonsingular projective variety $X$ over a field of characteristic $0$ and every $p$, there is a finite decreasing filtration $F^\bullet$ on $\CH^p(X)_\QQ$ such that:
\begin{enumerate}
\item $F^0 = \CH^p(X)_\QQ$, $F^1 = \CH^p(X)_{\mathrm{hom},\QQ}$ and $F^{p+1} = 0$;
\item the filtrations are compatible with intersection products, $F^i \cdot F^j \subseteq F^{i+j}$, and with the action of correspondences;
\item a correspondence $\Gamma \in \CH^{\dim X + r}(X \times Y)_\QQ$ acts on $\mathrm{Gr}^\nu_F\CH^p(X)_\QQ$ through its cohomology class, and it acts by $0$ if $\Gamma_* \colon H^{2p-\nu}(X) \to H^{2p-\nu+2r}(Y)$ vanishes.
\end{enumerate}
\end{conjecture}

See \cite{jannsen-1994} and \cite{bloch-lectures}; the formulation is due to Beilinson and Bloch, and Murre gave an equivalent one in terms of a decomposition of the diagonal into orthogonal idempotents
lifting the K\"unneth components to Chow groups \cite{murre-1993}. On $\mathrm{Gr}^1_F$ the filtration is expected to be detected by the Abel--Jacobi map (Section~\ref{vhs:section-intermediate-jacobians}), and
$\mathrm{Gr}^\nu_F$ is expected to be an $\mathrm{Ext}^\nu$ in a category of mixed motives. Beilinson and Bloch conjectured moreover that for $X$ defined over a number field, $F^2\CH^p(X_{\bar{\QQ}})_\QQ = 0$,
so that the Abel--Jacobi map is injective on $\CH^p(X_{\bar{\QQ}})_{\mathrm{hom},\QQ}$ \cite{beilinson-1987}.

\begin{remark}\label{absolute-hodge:remark-bb-bloch}
Conjecture~\ref{absolute-hodge:conjecture-bloch-beilinson} gives Bloch's conjecture up to torsion. Let $S$ be a surface with $h^{2,0}(S) = 0$. By the Lefschetz $(1,1)$-theorem, $H^2(S; \QQ)$ is spanned by
classes of divisors, so there is a correspondence $\Gamma = \sum a_{ij}D_i \times D_j$, with divisors $D_i$ and $a_{ij} \in \QQ$, that acts on $H^2(S)$ as the identity. On $0$-cycles $\Gamma_* = 0$, since
$(D_i \times D_j) \cdot (z \times S) = 0$ for a $0$-cycle $z$ in general position. The diagonal also acts on $H^2$ as the identity, so by (3) with $p = \nu = 2$, the identity of $\mathrm{Gr}^2_F\CH_0(S)_\QQ$ equals $\Gamma_* = 0$, and
$F^2\CH_0(S)_\QQ = 0$. The Albanese kernel is expected to lie in $F^2$, which is the case when the K\"unneth component $\pi_3$ is realised, as Murre showed it can be, by a correspondence acting through the
Albanese variety \cite{murre-1993}; then the Albanese kernel is torsion. Kimura's notion of \emph{finite-dimensional} motives gives another route: motives of curves and abelian varieties are
finite-dimensional, and finite-dimensionality of the motive of $S$ would imply Bloch's conjecture for $S$ \cite{kimura-2005}.
\end{remark}

\begin{remark}\label{absolute-hodge:remark-two-halves}
The Hodge conjecture describes $\mathrm{Gr}^0_F$: the image of the cycle class map is as large as Hodge theory permits. The Bloch--Beilinson conjectures describe $F^1$: the kernel is as structured as
motivic expectations require. Neither is known to imply the other, but a complete theory of algebraic cycles needs both, and both are expected to come from a category of mixed motives. The absolute Hodge
classes of Chapter~\ref{absolute-hodge}, the motivated cycles of Andr\'e and the unconditional results of these sections all live on the cohomological side; on the side of cycles, the Lefschetz $(1,1)$-theorem
(Section~\ref{hodge-conjecture:section-lefschetz-11}) remains the only general method for producing an algebraic cycle from a cohomology class.
\end{remark}

\section{Exercises on The Standard Conjectures and Motives}\label{absolute-hodge:section-exercises-standard-conjectures}

\begin{exercise}\label{absolute-hodge:exercise-kunneth-hodge}
Let $X$ be a nonsingular complex projective variety. Show that the K\"unneth components $\pi_i$ of the class of the diagonal are Hodge classes on $X \times X$, and deduce $C(X)$ from the Hodge conjecture
for $X \times X$ without using $B(X)$.
\end{exercise}

\begin{solution}
The class of the diagonal is algebraic, hence a Hodge class. The K\"unneth decomposition $H^{2n}(X \times X) = \bigoplus_iH^{2n-i}(X) \otimes H^i(X)$ is a decomposition into sub-Hodge structures, and the components
of a Hodge class along such a decomposition are Hodge classes (the components of a rational class of type $(n,n)$ are rational and of type $(n,n)$). The Hodge conjecture for $X \times X$ makes each $\pi_i$
algebraic.
\end{solution}

\begin{exercise}\label{absolute-hodge:exercise-projective-space}
Show that $h(\PP^n) \cong \bigoplus_{i=0}^n\mathbb{L}^{\otimes i}$ in $\mathcal{M}_{\mathrm{rat}}(k)$, and that $B(\PP^n)$ holds, by writing $\Lambda$ explicitly as a correspondence.
\end{exercise}

\begin{solution}
Let $h_1, h_2$ be the pullbacks of the hyperplane class to $\PP^n \times \PP^n$ and $p_i = h_1^{n-i}h_2^i \in \CH^n(\PP^n \times \PP^n)$. Then $(p_i)_*(h^j) = \mathrm{pr}_{2*}(h_1^{j+n-i}h_2^i) = \deg(h^{j+n-i})h^i = \delta_{ij}h^i$, and
$p_a \circ p_b = \mathrm{pr}_{13*}(h_1^{n-b}h_2^{b}h_2^{n-a}h_3^a) = \deg(h^{n+b-a})h_1^{n-b}h_3^a = \delta_{ab}p_a$. The group $\CH^n(\PP^n \times \PP^n)$ is free on the classes $h_1^{n-i}h_2^i$, and
$\gamma = \sum_ia_ih_1^{n-i}h_2^i$ acts by $\gamma_*(h^j) = a_jh^j$; as $\Delta$ acts as the identity, $\Delta = \sum_ip_i$. The correspondences $a = h^i \in \CH^i(\Spec k \times \PP^n)$, a morphism $\mathbb{L}^{\otimes i} \to h(\PP^n)$,
and $b = h^{n-i} \in \CH^{n-i}(\PP^n \times \Spec k)$, a morphism $h(\PP^n) \to \mathbb{L}^{\otimes i}$, satisfy $b \circ a = \deg(h^n) = 1$ and $a \circ b = h_1^{n-i}h_2^i = p_i$, so $(\PP^n, p_i, 0) \cong \mathbb{L}^{\otimes i}$. For $B$: the
only primitive class is $1 \in P^0$, so $\Lambda(h^j) = j(n - j + 1)h^{j-1}$. The correspondence $\lambda = \sum_{i=1}^ni(n - i + 1)h_1^{n-i}h_2^{i-1}$ satisfies
$\lambda_*(h^j) = \sum_ii(n - i + 1)\deg(h^{j+n-i})h^{i-1} = j(n - j + 1)h^{j-1}$, so $\lambda_* = \Lambda$ and $\Lambda$ is algebraic.
\end{solution}

\begin{exercise}\label{absolute-hodge:exercise-hypersurface}
Let $X \subseteq \PP^{n+1}$ be a nonsingular hypersurface over $\CC$. Show that $B(X)$ holds.
\end{exercise}

\begin{solution}
By the Lefschetz hyperplane theorem and Poincar\'e duality, $H^i(X; \QQ)$ is $\QQ h^{i/2}$ for $i \neq n$ even, $0$ for $i \neq n$ odd, and $H^n = \QQ h^{n/2} \oplus H^n_{\mathrm{prim}}$ (with the first summand only for $n$ even),
where $H^n_{\mathrm{prim}} = P^n$. The operator $\Lambda$ kills $P^n$ and acts on the powers of $h$ by $\Lambda(h^i) = i(n - i + 1)h^{i-1}$. Let $d = \deg X = \int h^n$, and consider the algebraic class
$\lambda = \frac1d\sum_{j=1}^nj(n - j + 1)h_1^{n-j}h_2^{j-1} \in H^{2n-2}(X \times X)$, where $h_1$, $h_2$ are the pullbacks of $h$. It acts by $\lambda_*(x) = \frac1d\sum_jj(n - j + 1)\big(\int_Xx \cup h^{n-j}\big)h^{j-1}$, and for
$x \in H^i$ only the term with $i = 2j$ can contribute. On $h^j$ it gives $j(n - j + 1)h^{j-1} = \Lambda(h^j)$, because $\int h^n = d$. For $x \in P^n$, the only possible term has $j = n/2$ (so $n$ is even), and
$\int x \cup h^{n/2} = \int L^{n/2}x = 0$ because $Lx = 0$ by primitivity (if $n \geq 2$; the case $n = 0$ is trivial). So $\lambda_*$ and $\Lambda$ agree on a basis, and $\Lambda$ is algebraic. The same argument applies to
smooth complete intersections, whose cohomology has the same shape.
\end{solution}

\begin{exercise}\label{absolute-hodge:exercise-jannsen-not-d}
Explain why Jannsen's theorem does not imply the standard conjecture $D$, although both concern numerical equivalence. Show, using the converse part of Theorem~\ref{absolute-hodge:theorem-jannsen},
that $D(X)$ for all $X$ is equivalent to the semisimplicity of $\mathcal{M}_{\mathrm{hom}}(k)$.
\end{exercise}

\begin{solution}
Jannsen's theorem is a statement about the quotient $A_{\mathrm{num}}$ of $A_{\mathrm{hom}}$; $D$ says that the kernel of $A_{\mathrm{hom}} \to A_{\mathrm{num}}$ vanishes. A quotient algebra can be semisimple whatever the kernel is, so
the first says nothing about the second. If $D$ holds for all varieties, then $\mathcal{M}_{\mathrm{hom}}(k) = \mathcal{M}_{\mathrm{num}}(k)$, which is semisimple abelian by Jannsen. Conversely, if
$\mathcal{M}_{\mathrm{hom}}(k)$ is semisimple abelian, the converse part of the theorem says that homological equivalence is numerical equivalence, which is $D$ for all varieties.
\end{solution}

\begin{exercise}\label{absolute-hodge:exercise-chow-plane}
Show that $\CH_0(\PP^2) = \ZZ$, via the degree, and $\CH_0(\PP^1 \times \PP^1) = \ZZ$, in accordance with Bloch's conjecture.
\end{exercise}

\begin{solution}
Two points of $\PP^2$ lie on a line $\ell \cong \PP^1$, and any two points of $\PP^1$ are rationally equivalent, being the zero and the pole of a rational function; pushing forward along $\ell \to \PP^2$, the two points are
rationally equivalent on $\PP^2$. So $\CH_0(\PP^2)$ is generated by the class of one point, and the degree shows it is $\ZZ$. On $\PP^1 \times \PP^1$, the points $(a, b)$ and $(a', b')$ are joined by the chain
$(a, b), (a', b), (a', b')$, whose consecutive points lie on a fibre $\cong \PP^1$. Both surfaces have $h^{2,0} = 0$ and trivial Albanese, so this is Bloch's conjecture for them.
\end{solution}
